\section{Case-split bookkeeping: the threefold cover and the central 34-point split}\label{cb:sec}

\subsection{Overview}\label{cb:sec:overview}

This section is about two Lean files that check the \emph{bookkeeping} of two case splits of the project, that is, the
purely combinatorial statement that an explicitly named family of nodes covers the split.

\begin{enumerate}[label=(\arabic*)]
\item The file \dirn{C3LEAN/C3Cover.lean} models the split of threefold rotationally symmetric sets of 33 points
  (11 orbits of the rotation by $120^\circ$, indexed by the residues $0,\dots,10$) by their convex layers; such sets
  are the subject of the case class \lean{Threefold33} of Definition~\ref{sr:def:classes}. Its headline is
  Theorem~\ref{cb:thm:c3} (\lean{C3Cover.all_roots_covered}): each of the 910 depth-two \emph{roots} of the reduced
  split is \emph{covered}, through layer splits and $\sigma$-splits, by the nodes named by the keys of 910 generated
  certificates. The proof combines a soundness lemma for a certificate checker (Lemma~\ref{cb:lem:run}, with a
  complete written proof here) with a kernel evaluation of the checker on the 910 certificates
  (Fact~\ref{cb:fact:batches}), the computer-assisted match of the codes of the roots with the roots of the
  certificates (Fact~\ref{cb:fact:roots-codes}), and the decoding of each root from its code
  (Fact~\ref{cb:fact:roots-match}, which also has a written proof here).
\item The file \dirn{C2FINAL/Central34Orbits.lean} models the second-layer split of the parallelogram hull cases of
  centrally symmetric sets of 34 points (17 antipodal classes $0,\dots,16$; the case class \lean{HalfTurn34} of
  Definition~\ref{sr:def:classes}). Its headline is Theorem~\ref{cb:thm:c2} (\lean{Central34.all_cases_closed}): for
  each hull case $P_a$, $1\le a\le 8$, each of the 560 second-layer subcases, or its image under the reflection that
  sends $i$ to the residue of $a-i$ modulo $17$, belongs (as a set) to an explicit generated list.
\end{enumerate}

\paragraph{What is not established here.} Both headline theorems are statements about finite lists of natural numbers.
They mention no points, no convexity, no CNF formula and no SAT solver, and they take no hypotheses. The following are
\emph{not} established by them, and are not proved in this document:
\begin{enumerate}[label=(\roman*)]
\item the geometric meaning of the splits: that every set of the class lies below some root (the top cases, and the
  reduction of the second layer by Lemma~5 and Lemma~H6M of the project's written arguments \dirn{C3-S/soundness.md}
  and \dirn{C3H6M/soundness_h6m.md}; for the central split, the reductions A and B reviewed in \dirn{C2P/review}),
  and that a layer split or a $\sigma$-split loses no set (Lemma~4 of \dirn{C3-S/soundness.md}, rules~2 and~3 of
  \dirn{C3U2/sigma_split.md}; the second-layer covers of the central split);
\item the CNF encodings of the nodes;
\item the SAT verdicts of the leaves, and even the fact that the leaves are nodes for which the project holds a run
  record in which both proof checkers, \texttt{cake\_lpr} and \texttt{lrat-check}, are recorded as accepting the proof
  (the project's term for such a record: ``certified''). In Theorem~\ref{cb:thm:c3} the set of leaves is, by
  definition, the set of nodes named by the keys of the certificates themselves, and in Theorem~\ref{cb:thm:c2} the
  lists $\mathsf{closedP}_a$ are just
  lists. Which nodes were put there was decided outside Lean from the run records (by
  \dirn{C3LEAN/tools/gen_c3lean.py}, audited by \dirn{C3LEAN/tools/audit_keys.py} and
  \dirn{C3LEAN/review/indep_audit.py}; by \dirn{C2FINAL/gen_lean.py} with \dirn{C2FINAL/hashmatch.py}).
\end{enumerate}
So neither theorem implies \lean{Threefold33} or \lean{HalfTurn34}. They check one step, ``the named nodes form a cut
of the split'', which the project combines with (i)--(iii); see Remarks~\ref{cb:rem:use} and~\ref{cb:rem:c2-scope}.

\paragraph{Sources and toolchain.} Paths are relative to the results directory \dirn{cloud/es7-2026-10-19/results}
of the repository.
\begin{itemize}
\item \dirn{C3LEAN/C3Cover.lean} (2{,}327 lines, SHA-256 \texttt{4b64293df2ea}\dots). Its lines 29--329 (the
  definitions of the model and of the checker, and the conceptual proofs) are byte-identical to lines 1--301 of
  \dirn{C3LEAN/tools/prelude.lean} (SHA-256 \texttt{cdcf96de7739}\dots), from which the generator
  \dirn{C3LEAN/tools/gen_c3lean.py} copies them. Lines 330--2299 hold the auxiliary definitions \lean{beqLayer},
  \lean{isOption}, \lean{distinctCodes}, \lean{codeSet} and \lean{reflectionFacts} (lines 335--365), the facts about
  the roots (lines 367--409 and 1720--1733), the generated data (lines 413, 423--1717 and 1737), the batch theorems
  (lines 1740--2245), the theorem \lean{all_checked} (line 2247), the main theorem (line 2251) and seven negative
  controls (lines 2266--2297); lines 2301--2327 are \texttt{\#print axioms} commands.
\item \dirn{C2FINAL/Central34Orbits.lean} (298 lines, SHA-256 \texttt{b396a26b2430}\dots), written by
  \dirn{C2FINAL/gen_lean.py}.
\end{itemize}
Both files are checked with Lean~4.24.0 according to their READMEs; for \dirn{C3LEAN/C3Cover.lean} a compile log is
saved, for \dirn{C2FINAL/Central34Orbits.lean} no Lean output (paragraph ``Axioms'' below). Neither has an \texttt{import}: only the core library is used, and the
repository's Mathlib (pinned at \texttt{f897ebcf}) plays no role. Every finite computation is closed by the tactic
\texttt{decide +kernel}: the Lean kernel itself evaluates the Boolean expression to \texttt{true} by definitional
unfolding, using its built-in arbitrary-precision arithmetic on natural-number literals (part of the kernel, not an
axiom). No \texttt{native\_decide} is used, so the axiom \lean{Lean.ofReduceBool} does not occur. The option
\texttt{Elab.async false} set at line~27 of \dirn{C3LEAN/C3Cover.lean} only affects elaboration.

\paragraph{Axioms.} For \dirn{C3LEAN/C3Cover.lean}, the compile log \dirn{C3LEAN/logs/compile.log} (and the identical
axiom lines of the review's recompilation, \dirn{C3LEAN/review/logs/main_compile.log}) records: \lean{all_roots_covered},
\lean{run_sound} and \lean{checkRoot_sound} depend on \lean{propext} and \lean{Quot.sound}; the eleven theorems
\lean{reflection_facts_H3}, \lean{reflection_facts_H6_1_0}, \dots, \lean{reflection_facts_H6_5_1}, and
\lean{roots_eq_tsv} and \lean{all_checked} on \lean{propext} only; \lean{topCases_eq}, \lean{strata},
\lean{roots_match} and \lean{roots_codes_match} on no axiom. The lemmas \lean{Covered.mono}, \lean{beqList_eq},
\lean{beqTok_eq}, \lean{beqNode_eq}, \lean{all_append_of} and the 127 batch theorems are not printed separately; each is
used by a printed theorem (\lean{run_sound}, \lean{all_checked} or \lean{all_roots_covered}), so its axioms are among
\lean{propext} and \lean{Quot.sound}. No theorem of the file depends on \lean{Classical.choice}, \lean{sorryAx} or
\lean{Lean.ofReduceBool}. For \dirn{C2FINAL/Central34Orbits.lean} no compile or axioms log is committed: the file ends
with three \texttt{\#print axioms} commands (lines 296--298) whose output is not in the repository, and
\dirn{C2FINAL/README.md} (\S6) records ``axioms \texttt{propext} only, about 6 s''. As no Lean output of any kind
(compile or build log, axiom output, pinned message) is saved for this file, the labels of its declarations rest on
that record alone: its three theorems and its definitions are \emph{Lean-checked according to the record}, the
qualified label defined with the rigor labels at the beginning of this document.

\paragraph{Rigor labels and dependencies.} The Lean theorems of \dirn{C3LEAN/C3Cover.lean} keep the label
\emph{Lean-checked}, and the three theorems of \dirn{C2FINAL/Central34Orbits.lean} (Theorem~\ref{cb:thm:c2-refl},
Fact~\ref{cb:fact:c2-sub} and Theorem~\ref{cb:thm:c2}) carry the label \emph{Lean-checked according to the record},
as do its definitions (paragraph ``Axioms'' above); the proofs written
here are \emph{written proofs (not refereed)}. Every statement that Lean proves by kernel evaluation is marked
\emph{computer-assisted} in its statement, and so is every statement whose proof uses one of the computer-assisted
statements on generated data (Lemma~\ref{cb:lem:all-checked}, Corollary~\ref{cb:cor:roots},
Theorem~\ref{cb:thm:c3} and Corollary~\ref{cb:cor:c2}). For the kernel evaluations on explicit generated data
(Facts~\ref{cb:fact:batches}, \ref{cb:fact:roots-codes}, \ref{cb:fact:roots-tsv} and~\ref{cb:fact:c2-sub}, and
Theorem~\ref{cb:thm:c2}) we state exactly what is evaluated, and we give no hand proof. For the kernel evaluations of
the definitions alone (Facts~\ref{cb:fact:refl-h3} to~\ref{cb:fact:roots-match} and Theorem~\ref{cb:thm:c2-refl}) we
also give complete written proofs that do not use the computation; the statement names the written proof. Each Lean
declaration of the two files other than the seven negative controls of \dirn{C3LEAN/C3Cover.lean} (tests, listed in
\S\ref{cb:sec:coverage}), or each uniform family of generated declarations (\lean{certs0}, \dots, \lean{certs126};
\lean{batch0}, \dots, \lean{batch126}; the ten theorems \lean{reflection_facts_H6_1_0}, \dots,
\lean{reflection_facts_H6_5_1}; \lean{closedP1}, \dots, \lean{closedP8}), has a statement below whose head carries
its Lean reference; for a family the head cites the first member, and the statement locates the others.
The head of Lemma~\ref{cb:lem:run} carries the two declarations \lean{run_sound} and \lean{checkRoot_sound}. The
proofs of this section use no other section.

\subsection{Notation}\label{cb:sec:notation}

\paragraph{Numbers.} $\N=\{0,1,2,\dots\}$. For $x,y\in\N$ with $y>0$, $\lfloor x/y\rfloor$ and $x\bmod y$ are the
quotient and the remainder of Euclidean division (Lean's division and remainder on $\N$), and
$x\mathbin{\dot-}y=\max(x-y,0)$ is Lean's truncated subtraction on $\N$. For an integer $x$ and $y\in\N$ with
$y>0$, $x\bmod y$ denotes the least nonnegative residue of $x$ modulo $y$; for $x\in\N$ this is the remainder above.
For $\sigma,f\in\{0,1\}$,
$\sigma\oplus f=(\sigma+f)\bmod 2$.

\paragraph{Lists.} We write $[x_0,\dots,x_{k-1}]$ for a list, $[]$ for the empty list, $x\mathbin{::}l$ for the list
with head $x$ and tail $l$, $l_1\mathbin{+\!\!+}l_2$ for concatenation, $\card{l}$ for the length, and $x\in l$ if $x$
is an entry of $l$; $\mathrm{set}(l)$ is the set of entries of $l$. A list is \emph{duplicate-free} if entries at
different positions are different; then $\card{l}=\card{\mathrm{set}(l)}$. The list functions are those of Lean's core
library: $\mathsf{range}(n)=[0,1,\dots,n-1]$; $\mathsf{map}(g,l)=[g(x_0),\dots,g(x_{k-1})]$;
$\mathsf{flatMap}(g,l)=g(x_0)\mathbin{+\!\!+}\cdots\mathbin{+\!\!+}g(x_{k-1})$; $\mathsf{filter}(p,l)$ is the list of
the entries $x$ of $l$ with $p(x)=\mathsf{true}$, in their order; $\mathsf{take}(K,l)$ is the list of the first
$\min(K,\card l)$ entries; $\mathsf{zip}(l,l')=[(x_0,x'_0),(x_1,x'_1),\dots]$ has length $\min(\card l,\card{l'})$;
$\mathsf{foldl}(g,z,l)=g(\cdots g(g(z,x_0),x_1)\cdots,x_{k-1})$; $\mathsf{replicate}(k,x)$ is the list of $k$ copies of
$x$. The Boolean-valued functions are: $\mathsf{contains}(l,x)=\mathsf{true}$ iff $x\in l$;
$\mathsf{all}(l,p)=\mathsf{true}$ iff $p(x)=\mathsf{true}$ for every $x\in l$; $\mathsf{any}(l,p)=\mathsf{true}$ iff
$p(x)=\mathsf{true}$ for some $x\in l$; $\mathsf{isEmpty}(l)=\mathsf{true}$ iff $l=[]$. We use without comment the
following facts, each immediate by induction on $l$ from the defining equations:
\begin{enumerate}[label=(L\arabic*)]
\item $\mathsf{all}(l_1\mathbin{+\!\!+}l_2,p)=\mathsf{all}(l_1,p)\wedge\mathsf{all}(l_2,p)$;
\item $y\in\mathsf{flatMap}(g,l)$ iff $y\in g(x)$ for some $x\in l$;
\item $\card{\mathsf{map}(g,l)}=\card l$, and $\card{\mathsf{flatMap}(g,l)}$ is the sum of $\card{g(x)}$ over the
  positions of $l$;
\item $y\in\mathsf{map}(g,l)$ iff $y=g(x)$ for some $x\in l$.
\end{enumerate}

\paragraph{Booleans.} The Booleans are $\mathsf{true}$ and $\mathsf{false}$. In Boolean expressions $\wedge$, $\vee$,
$\neg$ denote Lean's Boolean conjunction, disjunction and negation, and a decidable proposition $P$ used where a Boolean
is expected stands for $\mathsf{decide}(P)$, the Boolean that is $\mathsf{true}$ iff $P$ holds (Lean inserts this
coercion). For $x,y\in\N$, Lean's Boolean equality tests (written \texttt{x == y} or \texttt{Nat.beq x y}) are
$\mathsf{true}$ iff $x=y$, and \texttt{x != y} is the negation of \texttt{x == y}. A statement ``$b=\mathsf{true}$'' is a
Lean proposition about a Boolean $b$.

\paragraph{Binary digits.} For $n\in\N$ let $\beta(n)=\{i\in\N:\lfloor n/2^i\rfloor\text{ is odd}\}$, the set of
positions of the digit $1$ in the binary expansion of $n$. Then $n=\sum_{i\in\beta(n)}2^i$, and $n\mapsto\beta(n)$ is a
bijection from $\N$ onto the set of finite subsets of $\N$ (uniqueness of the binary expansion). In Lean, shifting $1$
left by $c$ places gives $2^c$, so $\beta(2^c)=\{c\}$; the bitwise or of $x$ and $y$, written here
$x\mathbin{\mathsf{or}}y$, is the number with $\beta(x\mathbin{\mathsf{or}}y)=\beta(x)\cup\beta(y)$; and
$\mathsf{testBit}(n,i)$ (Lean's \texttt{Nat.testBit n i}) is $\mathsf{true}$ iff $i\in\beta(n)$.

\begin{lemma}[Or-folds; written step]\label{cb:lem:bits}
Let $z\in\N$, let $l$ be a list of elements of a set $X$, and let $g:X\to\N$. Then
\[
\beta\big(\mathsf{foldl}((acc,x)\mapsto acc\mathbin{\mathsf{or}}2^{g(x)},\ z,\ l)\big)=\beta(z)\cup\{g(x):x\in l\}.
\]
\end{lemma}

\begin{proof}
Induction on $l$, for all $z$ at once. For $l=[]$ the fold is $z$. For $l=x\mathbin{::}l'$ the fold equals
$\mathsf{foldl}((acc,y)\mapsto acc\mathbin{\mathsf{or}}2^{g(y)},\ z\mathbin{\mathsf{or}}2^{g(x)},\ l')$, whose set of
binary digits is, by induction, $\beta(z)\cup\{g(x)\}\cup\{g(y):y\in l'\}$.
\end{proof}

\subsection{The threefold split: the model}\label{cb:sec:model}

Every Lean reference in a statement head of \S\S\ref{cb:sec:model}--\ref{cb:sec:c3} is to
\dirn{C3LEAN/C3Cover.lean}, namespace \lean{C3Cover}; the statements tagged ``written step'' have no Lean
counterpart. The definitions transcribe the project's Python implementation of the split
(\dirn{C3U2/tools/hullcase.py}); that the transcription agrees with the Python code is tested in the project
(\dirn{C3LEAN/README.md}, \S3), not proved, and plays no role in the theorems.

\begin{definition}[Orbits; \leanref{C3Cover.m}{C3LEAN/C3Cover.lean:34}]\label{cb:def:m}
$m=11$. Below we write $11$ for $m$ wherever it occurs.
\end{definition}

\begin{definition}[Tokens; \leanref{C3Cover.Tok}{C3LEAN/C3Cover.lean:38}]\label{cb:def:tok}
A \emph{token} is a pair $t=(r,s)\in\N\times\N$; we write $t.1=r$ and $t.2=s$.
\end{definition}

\begin{definition}[Layers; \leanref{C3Cover.Layer}{C3LEAN/C3Cover.lean:41}]\label{cb:def:layer}
A \emph{layer} is a list of tokens.
\end{definition}

\begin{definition}[Nodes; \leanref{C3Cover.Node}{C3LEAN/C3Cover.lean:45}]\label{cb:def:node}
A \emph{node} $n=\langle\Lambda,C\rangle$ is a pair consisting of a list $n.\mathsf{layers}=\Lambda$ of layers and a
list $n.\mathsf{cube}=C$ of pairs $(r,v)\in\N\times\N$. Two nodes are equal iff both components are equal.
\end{definition}

\begin{definition}[Used residues; \leanref{C3Cover.Node.used}{C3LEAN/C3Cover.lean:51}]\label{cb:def:used}
For a node $n$,
\[
\mathsf{used}(n)=\mathsf{flatMap}\big(L\mapsto\mathsf{map}(t\mapsto t.1,\ L),\ n.\mathsf{layers}\big),
\]
the first components of the tokens of all layers, in order.
\end{definition}

\begin{remark}[Intended reading]\label{cb:rem:reading}
In the project, a token $(r,s)$ stands for the orbit with residue $r\in\{0,\dots,10\}$ together with its $\sigma$
value $s\in\{0,1\}$ (the side of a fixed line through the centre on which the point with label $r$ lies), and $s=2$
stands for ``no $\sigma$'', which is used for a layer of one orbit other than the top layer (the one-orbit top
layer $[(0,0)]$ carries $s=0$). A layer of one or two tokens is a convex layer of
one or two orbits; the layers of a node are listed from the convex hull inwards, and the cube fixes the $\sigma$ values
$v$ of some further residues $r$. The top layer is $\mathsf{H3}=[(0,0)]$ (a hull of one orbit) or
$\mathsf{H6}(a,s)=[(0,0),(a,s)]$ with $1\le a\le5$ and $s\in\{0,1\}$ (a hull of the orbits $0$ and $a$, with
$\sigma_a=s$); see Lemma~2 of \dirn{C3-S/soundness.md}. Nothing in the Lean types restricts $r\le10$ or $s\le2$: all such
bounds used below are consequences of how particular lists are built. This reading is not used in any proof.
\end{remark}

\begin{definition}[Pair options; \leanref{C3Cover.pairOptions}{C3LEAN/C3Cover.lean:54}]\label{cb:def:pairoptions}
For a list $l$ of naturals,
\[
\mathsf{pairOptions}([])=[],\qquad
\mathsf{pairOptions}(b\mathbin{::}l)=\mathsf{flatMap}\big(c\mapsto Q_{b,c},\ l\big)\mathbin{+\!\!+}\mathsf{pairOptions}(l),
\]
where $Q_{b,c}=\big[\,[(b,0),(c,0)],\ [(b,0),(c,1)],\ [(b,1),(c,0)],\ [(b,1),(c,1)]\,\big]$.
\end{definition}

\begin{definition}[Layer options; \leanref{C3Cover.layerOptions}{C3LEAN/C3Cover.lean:61}]\label{cb:def:layeroptions}
For a list $U$ of naturals let $\mathsf{rem}(U)=\mathsf{filter}\big(r\mapsto\neg\mathsf{contains}(U,r),\
\mathsf{range}(11)\big)$ (a local abbreviation inside \lean{layerOptions}), and
\[
\mathsf{layerOptions}(U)=\mathsf{map}\big(r\mapsto[(r,2)],\ \mathsf{rem}(U)\big)\mathbin{+\!\!+}
\mathsf{pairOptions}(\mathsf{rem}(U)).
\]
\end{definition}

\begin{definition}[Insertion; \leanref{C3Cover.insTok}{C3LEAN/C3Cover.lean:66}]\label{cb:def:instok}
For a pair $t\in\N\times\N$ and a list of such pairs: $\mathsf{insTok}(t,[])=[t]$;
$\mathsf{insTok}(t,u\mathbin{::}us)=t\mathbin{::}u\mathbin{::}us$ if $t.1\le u.1$, and
$\mathsf{insTok}(t,u\mathbin{::}us)=u\mathbin{::}\mathsf{insTok}(t,us)$ if $t.1>u.1$.
\end{definition}

\begin{definition}[Sorting; \leanref{C3Cover.sortTok}{C3LEAN/C3Cover.lean:71}]\label{cb:def:sorttok}
$\mathsf{sortTok}([])=[]$ and $\mathsf{sortTok}(t\mathbin{::}ts)=\mathsf{insTok}(t,\mathsf{sortTok}(ts))$.
\end{definition}

\begin{definition}[Top cases; \leanref{C3Cover.topCases}{C3LEAN/C3Cover.lean:77}]\label{cb:def:topcases}
The list of top cases is
\[
\mathsf{topCases}=[(0,0)]\mathbin{::}\mathsf{flatMap}\big(i\mapsto\big[\,[(0,0),(i+1,0)],\ [(0,0),(i+1,1)]\,\big],\
\mathsf{range}(5)\big).
\]
We write $\mathsf{H3}=[(0,0)]$ and $\mathsf{H6}(a,s)=[(0,0),(a,s)]$.
\end{definition}

\begin{definition}[Option keys; \leanref{C3Cover.okey}{C3LEAN/C3Cover.lean:82}]\label{cb:def:okey}
For a layer $L$, $\mathsf{okey}(L)=\mathsf{map}(\kappa,\mathsf{sortTok}(L))$, where $\kappa(r,s)=3r$ if $s=2$ and
$\kappa(r,s)=3r+s+1$ if $s\ne2$.
\end{definition}

\begin{definition}[Lexicographic order; \leanref{C3Cover.lexLe}{C3LEAN/C3Cover.lean:85}]\label{cb:def:lexle}
For lists of naturals: $\mathsf{lexLe}([],y)=\mathsf{true}$; $\mathsf{lexLe}(x\mathbin{::}xs,[])=\mathsf{false}$; and
$\mathsf{lexLe}(x\mathbin{::}xs,\,y\mathbin{::}ys)=(x<y)\vee\big((x=y)\wedge\mathsf{lexLe}(xs,ys)\big)$.
\end{definition}

\begin{definition}[Mirror; \leanref{C3Cover.mirror}{C3LEAN/C3Cover.lean:91}]\label{cb:def:mirror}
For a layer $L$,
\[
\mathsf{mirror}(L)=\mathsf{sortTok}(\mathsf{map}(\mu,L)),\qquad
\mu(r,s)=\big((11\mathbin{\dot-}r)\bmod11,\ s'\big),\quad s'=\begin{cases}2&\text{if }s=2,\\
1\mathbin{\dot-}s&\text{if }s\ne2.\end{cases}
\]
\end{definition}

\begin{definition}[The map phi; \leanref{C3Cover.phi}{C3LEAN/C3Cover.lean:95}]\label{cb:def:phi}
For $a,s\in\N$ and a layer $L$,
\[
\mathsf{phi}(a,s,L)=\mathsf{sortTok}(\mathsf{map}(\varphi_{a,s},L)),\qquad
\varphi_{a,s}(r,\sigma)=\big(((a+11)\mathbin{\dot-}r)\bmod11,\ \sigma'\big),
\]
where $\sigma'=2$ if $\sigma=2$; $\sigma'=1\mathbin{\dot-}\sigma$ if $\sigma\ne2$ and exactly one of the conditions
$r<a$, $s=0$ holds; and $\sigma'=\sigma$ if $\sigma\neq2$ and both or neither of them hold.
\end{definition}

\begin{definition}[Reflection; \leanref{C3Cover.refl}{C3LEAN/C3Cover.lean:101}]\label{cb:def:refl}
For layers $top$ and $o$: $\mathsf{refl}(top,o)=\mathsf{mirror}(o)$ if $top=[(0,0)]$;
$\mathsf{refl}(top,o)=\mathsf{phi}(a,s,o)$ if $top=[(0,0),(a,s)]$ for some $a,s\in\N$; and
$\mathsf{refl}(top,o)=o$ for every other list $top$.
\end{definition}

\begin{definition}[Kept options; \leanref{C3Cover.keptOptions}{C3LEAN/C3Cover.lean:109}]\label{cb:def:kept}
For a layer $top$,
\begin{align*}
\mathsf{keptOptions}(top)&=\mathsf{filter}\big(\mathsf{kept}_{top},\ \mathsf{layerOptions}(\mathsf{map}(u\mapsto
u.1,\ top))\big),\quad\text{where}\\
\mathsf{kept}_{top}(o)&=\mathsf{lexLe}\big(\mathsf{okey}(o),\ \mathsf{okey}(\mathsf{refl}(top,o))\big).
\end{align*}
\end{definition}

\begin{definition}[Roots; \leanref{C3Cover.roots}{C3LEAN/C3Cover.lean:113}]\label{cb:def:roots}
The list of roots is
\[
\mathsf{roots}=\mathsf{flatMap}\big(t\mapsto\mathsf{map}(o\mapsto\langle[t,o],[]\rangle,\ \mathsf{keptOptions}(t)),\
\mathsf{topCases}\big).
\]
\end{definition}

\begin{definition}[Free residues; \leanref{C3Cover.Node.free}{C3LEAN/C3Cover.lean:116}]\label{cb:def:free}
For a node $n$,
\[
\mathsf{free}(n)=\mathsf{filter}\big(r\mapsto(r\ne0)\wedge\neg\mathsf{contains}(\mathsf{used}(n),r)\wedge
\neg\mathsf{contains}(\mathsf{map}(c\mapsto c.1,\ n.\mathsf{cube}),r),\ \mathsf{range}(11)\big).
\]
\end{definition}

\begin{definition}[Bits; \leanref{C3Cover.bitsOf}{C3LEAN/C3Cover.lean:120}]\label{cb:def:bitsof}
For $K,i\in\N$,
\[
\mathsf{bitsOf}(K,i)=\mathsf{map}\big(j\mapsto\lfloor i/2^{(K\mathbin{\dot-}1)\mathbin{\dot-}j}\rfloor\bmod2,\
\mathsf{range}(K)\big).
\]
\end{definition}

\begin{definition}[Adding bits; \leanref{C3Cover.addBits}{C3LEAN/C3Cover.lean:123}]\label{cb:def:addbits}
For lists $C$ and $new$ of pairs of naturals,
\[
\mathsf{addBits}(C,new)=\mathsf{foldl}\big((c,t)\mapsto\mathsf{insTok}(t,c),\ C,\ new\big).
\]
\end{definition}

\begin{definition}[Sigma children; \leanref{C3Cover.Node.sigmaChildren}{C3LEAN/C3Cover.lean:127}]\label{cb:def:sigmachildren}
For a node $n$ and $K\in\N$, with the local abbreviation $bits=\mathsf{take}(K,\mathsf{free}(n))$,
\[
\mathsf{sigmaChildren}(n,K)=\mathsf{map}\Big(i\mapsto\big\langle n.\mathsf{layers},\
\mathsf{addBits}\big(n.\mathsf{cube},\ \mathsf{zip}(bits,\mathsf{bitsOf}(K,i))\big)\big\rangle,\
\mathsf{range}(2^K)\Big).
\]
\end{definition}

\begin{definition}[Conflicts; \leanref{C3Cover.conflicts}{C3LEAN/C3Cover.lean:132}]\label{cb:def:conflicts}
For a layer $L$ and a list $C$ of pairs of naturals,
\[
\mathsf{conflicts}(L,C)=\mathsf{any}\big(L,\ t\mapsto(t.2\ne2)\wedge\mathsf{any}(C,\ c\mapsto(c.1=t.1)\wedge(c.2\ne
t.2))\big).
\]
\end{definition}

\begin{definition}[Layer children; \leanref{C3Cover.Node.layerChildren}{C3LEAN/C3Cover.lean:137}]\label{cb:def:layerchildren}
For a node $n$,
\[
\mathsf{layerChildren}(n)=\mathsf{map}\Big(o\mapsto\langle n.\mathsf{layers}\mathbin{+\!\!+}[o],\
n.\mathsf{cube}\rangle,\ \mathsf{filter}\big(o\mapsto\neg\mathsf{conflicts}(o,n.\mathsf{cube}),\
\mathsf{layerOptions}(\mathsf{used}(n))\big)\Big).
\]
\end{definition}

\begin{remark}[Reading of the children]\label{cb:rem:children}
If $K\le\card{\mathsf{free}(n)}$ and $b_0<\dots<b_{K-1}$ are the first $K$ entries of $\mathsf{free}(n)$ (the $K$
smallest residues in $\{1,\dots,10\}$ that occur neither in a layer nor in the cube of $n$), then the $i$-th
$\sigma$-child, $0\le i<2^K$, has the layers of $n$ and the cube of $n$ with the pairs $(b_j,v_j)$ inserted, where
$v_j=\lfloor i/2^{K-1-j}\rfloor\bmod2$ is the binary digit of $i$ at position $K-1-j$; as $i$ runs through
$0,\dots,2^K-1$, $(v_0,\dots,v_{K-1})$ runs through $\{0,1\}^K$. A layer child of $n$ appends to the layers of $n$ one
layer option on the residues not used by $n$, except the options that give some residue of the cube an explicit
$\sigma$ value (not $2$) different from the cube's value. No reduction by a reflection is applied to layer children.
This reading is not used in any proof.
\end{remark}

\subsection{The cover predicate}\label{cb:sec:covered}

\begin{definition}[Covered nodes; \leanref{C3Cover.Covered}{C3LEAN/C3Cover.lean:146}]\label{cb:def:covered}
Let $S$ be a predicate on nodes. A predicate $P$ on nodes is \emph{$S$-closed} if for every node $n$ the following
three rules hold.
\begin{description}
\item[(cert)] If $S(n)$, then $P(n)$.
\item[(layer)] If $2\le\card{n.\mathsf{layers}}$, $\mathsf{layerChildren}(n)\ne[]$, and $P(c)$ for every
  $c\in\mathsf{layerChildren}(n)$, then $P(n)$.
\item[(sigma)] If $K\in\N$, $2\le\card{n.\mathsf{layers}}$, $0<K$, $K\le\card{\mathsf{free}(n)}$, and $P(c)$ for every
  $c\in\mathsf{sigmaChildren}(n,K)$, then $P(n)$.
\end{description}
$\mathsf{Covered}(S)$ is the least $S$-closed predicate: $\mathsf{Covered}(S)(n)$ holds iff $P(n)$ holds for every
$S$-closed predicate $P$.
\end{definition}

In each rule the predicate occurs in the premises only positively, so an intersection of $S$-closed predicates is
$S$-closed; hence $\mathsf{Covered}(S)$ is itself $S$-closed. We use the two resulting principles:
\emph{closure} ($\mathsf{Covered}(S)$ satisfies (cert), (layer), (sigma)) and \emph{induction} (if $P$ is $S$-closed,
then $\mathsf{Covered}(S)(n)$ implies $P(n)$ for every node $n$). In Lean, \texttt{Covered S} is an inductive predicate
whose three constructors \lean{Covered.cert}, \lean{Covered.layer}, \lean{Covered.sigma} are exactly these three rules,
with these premises in this order, and whose recursor yields the induction principle. The Lean predicate is therefore
$S$-closed and contained in every $S$-closed predicate, so it coincides with the predicate defined above.

\begin{lemma}[Monotonicity; \leanref{C3Cover.Covered.mono}{C3LEAN/C3Cover.lean:153}]\label{cb:lem:mono}
Let $S,T$ be predicates on nodes such that $S(n)$ implies $T(n)$ for every node $n$. Then for every node $n$,
$\mathsf{Covered}(S)(n)$ implies $\mathsf{Covered}(T)(n)$.
\end{lemma}

\begin{proof}
As in Lean, by induction on the derivation. The predicate $P=\mathsf{Covered}(T)$ is $S$-closed: for (cert), $S(n)$
gives $T(n)$, hence $P(n)$ by the rule (cert) for $T$; the rules (layer) and (sigma) have the same side conditions for
$S$ and for $T$, so they hold for $P$ because $\mathsf{Covered}(T)$ satisfies the rules (layer) and (sigma) for $T$. By
the induction principle, $\mathsf{Covered}(S)(n)$ implies $P(n)$.
\end{proof}

\subsection{Keys, codes and Boolean equality}\label{cb:sec:keys}

\begin{definition}[Token decoding; \leanref{C3Cover.decodeTok}{C3LEAN/C3Cover.lean:165}]\label{cb:def:decodetok}
For $x\in\N$, $\mathsf{decodeTok}(x)=(\lfloor x/3\rfloor,\ x\bmod3)$.
\end{definition}

\begin{definition}[Layer decoding; \leanref{C3Cover.decodeLayer}{C3LEAN/C3Cover.lean:167}]\label{cb:def:decodelayer}
For $c\in\N$,
\[
\mathsf{decodeLayer}(c)=\begin{cases}[\mathsf{decodeTok}(c\mathbin{\dot-}1)]&\text{if }c\le33,\\
\big[\mathsf{decodeTok}(\lfloor(c\mathbin{\dot-}34)/33\rfloor),\ \mathsf{decodeTok}((c\mathbin{\dot-}34)\bmod33)\big]
&\text{if }c>33.\end{cases}
\]
\end{definition}

\begin{definition}[Decoding layers; \leanref{C3Cover.decodeLayers}{C3LEAN/C3Cover.lean:172}]\label{cb:def:decodelayers}
For $f,lk\in\N$ and a list $acc$ of layers: $\mathsf{decodeLayers}(0,lk,acc)=acc$;
$\mathsf{decodeLayers}(f+1,lk,acc)=acc$ if $lk=0$; and $\mathsf{decodeLayers}(f+1,lk,acc)=\mathsf{decodeLayers}\big(f,\
\lfloor lk/2048\rfloor,\ \mathsf{decodeLayer}(lk\bmod2048)\mathbin{::}acc\big)$ if $lk\ne0$.
\end{definition}

\begin{definition}[Cube decoding; \leanref{C3Cover.decodeCube}{C3LEAN/C3Cover.lean:180}]\label{cb:def:decodecube}
For $f,r,cm,cv\in\N$: $\mathsf{decodeCube}(0,r,cm,cv)=[]$; $\mathsf{decodeCube}(f+1,r,cm,cv)=[]$ if $cm=0$; and if
$cm\neq0$, $\mathsf{decodeCube}(f+1,r,cm,cv)$ is $(r,\ cv\bmod2)\mathbin{::}D$ when $cm\bmod2=1$ and $D$ when
$cm\bmod2\ne1$, where $D=\mathsf{decodeCube}(f,\ r+1,\ \lfloor cm/2\rfloor,\ \lfloor cv/2\rfloor)$.
\end{definition}

\begin{definition}[Key decoding; \leanref{C3Cover.decodeKey}{C3LEAN/C3Cover.lean:191}]\label{cb:def:decodekey}
For $k\in\N$,
\[
\mathsf{decodeKey}(k)=\big\langle\mathsf{decodeLayers}(12,\ \lfloor k/2^{22}\rfloor,\ []),\
\mathsf{decodeCube}(11,\ 0,\ \lfloor k/2^{11}\rfloor\bmod2^{11},\ k\bmod2^{11})\big\rangle.
\]
(In Lean the constants are written $4194304=2^{22}$ and $2048=2^{11}$.)
\end{definition}

\begin{remark}[Keys]\label{cb:rem:keys}
Read in base $2048$, the number $\lfloor k/2^{22}\rfloor$ lists the layer codes of $\mathsf{decodeKey}(k)$, its most
significant digit giving the top layer (each step of $\mathsf{decodeLayers}$ puts the layer of the least significant
remaining digit in front of the layers decoded so far); the next $11$ binary digits of $k$ are the set of cube residues
and the last $11$ their values. For example (the review's theorem \lean{C3Cover.rev_parent_decodes},
\dirn{C3LEAN/review/RevCtrl_pass.lean:315}, whose Lean proof is \texttt{decide +kernel}; its file compiles with exit status~0,
\dirn{C3LEAN/review/logs/revctrl_pass.log}, but the axioms of this theorem are not printed),
$\mathsf{decodeKey}(8615616512)=\langle[[(0,0)],[(1,2)]],\ [(2,0),(3,0),(4,0),(5,0),(6,0),(7,0)]\rangle$. The map
$\mathsf{decodeKey}$ is not injective: value digits outside the residue mask are ignored, and a zero digit inside the
layer part decodes to the layer $[(0,0)]$. Nothing below needs injectivity: the leaves are defined as decoded keys, and
the checker compares nodes, not keys.
\end{remark}

\begin{definition}[Token code; \leanref{C3Cover.tokCode}{C3LEAN/C3Cover.lean:223}]\label{cb:def:tokcode}
For a token $t$, $\mathsf{tokCode}(t)=3\,t.1+t.2$.
\end{definition}

\begin{definition}[Layer code; \leanref{C3Cover.layerCode}{C3LEAN/C3Cover.lean:226}]\label{cb:def:layercode}
For tokens $t,u$: $\mathsf{layerCode}([t])=1+\mathsf{tokCode}(t)$,
$\mathsf{layerCode}([t,u])=34+33\,\mathsf{tokCode}(t)+\mathsf{tokCode}(u)$, and $\mathsf{layerCode}(L)=0$ for every
layer $L$ with no or more than two tokens.
\end{definition}

\begin{definition}[Root code; \leanref{C3Cover.rootCode}{C3LEAN/C3Cover.lean:232}]\label{cb:def:rootcode}
For a node $n$: $\mathsf{rootCode}(n)=\mathsf{layerCode}(A)\cdot1123+\mathsf{layerCode}(B)$ if
$n.\mathsf{layers}=[A,B]$, and $\mathsf{rootCode}(n)=0$ otherwise.
\end{definition}

\begin{definition}[Root decoding; \leanref{C3Cover.decodeRoot}{C3LEAN/C3Cover.lean:238}]\label{cb:def:decoderoot}
For $c\in\N$, $\mathsf{decodeRoot}(c)=\big\langle[\mathsf{decodeLayer}(\lfloor c/1123\rfloor),\
\mathsf{decodeLayer}(c\bmod1123)],\ []\big\rangle$.
\end{definition}

\begin{definition}[List equality; \leanref{C3Cover.beqList}{C3LEAN/C3Cover.lean:194}]\label{cb:def:beqlist}
For a type $\alpha$ and a Boolean-valued function $eq$ of two arguments in $\alpha$:
$\mathsf{beqList}(eq,[],[])=\mathsf{true}$;
$\mathsf{beqList}(eq,\,a\mathbin{::}as,\,b\mathbin{::}bs)=eq(a,b)\wedge\mathsf{beqList}(eq,as,bs)$; and
$\mathsf{beqList}(eq,l_1,l_2)=\mathsf{false}$ if exactly one of $l_1,l_2$ is empty.
\end{definition}

\begin{definition}[Token equality; \leanref{C3Cover.beqTok}{C3LEAN/C3Cover.lean:199}]\label{cb:def:beqtok}
For pairs $a,b$ of naturals, $\mathsf{beqTok}(a,b)=\mathsf{Nat.beq}(a.1,b.1)\wedge\mathsf{Nat.beq}(a.2,b.2)$.
\end{definition}

\begin{definition}[Node equality; \leanref{C3Cover.beqNode}{C3LEAN/C3Cover.lean:202}]\label{cb:def:beqnode}
For nodes $a,b$,
\[
\mathsf{beqNode}(a,b)=\mathsf{beqList}\big(\mathsf{beqList}(\mathsf{beqTok},\cdot,\cdot),\
a.\mathsf{layers},\ b.\mathsf{layers}\big)\wedge\mathsf{beqList}(\mathsf{beqTok},\ a.\mathsf{cube},\ b.\mathsf{cube}).
\]
\end{definition}

\begin{definition}[Layer equality; \leanref{C3Cover.beqLayer}{C3LEAN/C3Cover.lean:335}]\label{cb:def:beqlayer}
For layers $L,L'$, $\mathsf{beqLayer}(L,L')=\mathsf{beqList}(\mathsf{beqTok},L,L')$.
\end{definition}

\begin{lemma}[List equality; \leanref{C3Cover.beqList_eq}{C3LEAN/C3Cover.lean:204}]\label{cb:lem:beqlist}
Let $\alpha$ be a type and $eq$ a Boolean-valued function on $\alpha\times\alpha$ such that $eq(a,b)=\mathsf{true}$
implies $a=b$ for all $a,b\in\alpha$. Then for all lists $l_1,l_2$ over $\alpha$,
$\mathsf{beqList}(eq,l_1,l_2)=\mathsf{true}$ implies $l_1=l_2$.
\end{lemma}

\begin{proof}
As in Lean, by induction on $l_1$, for all $l_2$. If both lists are empty they are equal. If exactly one is empty,
$\mathsf{beqList}$ is $\mathsf{false}$, contrary to the assumption. If $l_1=a\mathbin{::}as$ and $l_2=b\mathbin{::}bs$,
the assumption gives $eq(a,b)=\mathsf{true}$ and $\mathsf{beqList}(eq,as,bs)=\mathsf{true}$; so $a=b$ by the hypothesis
on $eq$, and $as=bs$ by induction; hence $l_1=l_2$.
\end{proof}

\begin{lemma}[Token equality; \leanref{C3Cover.beqTok_eq}{C3LEAN/C3Cover.lean:213}]\label{cb:lem:beqtok}
For all tokens $a,b$: $\mathsf{beqTok}(a,b)=\mathsf{true}$ implies $a=b$.
\end{lemma}

\begin{proof}
$\mathsf{Nat.beq}(x,y)=\mathsf{true}$ implies $x=y$, so both components of $a$ and $b$ agree.
\end{proof}

\begin{lemma}[Node equality; \leanref{C3Cover.beqNode_eq}{C3LEAN/C3Cover.lean:218}]\label{cb:lem:beqnode}
For all nodes $a,b$: $\mathsf{beqNode}(a,b)=\mathsf{true}$ implies $a=b$.
\end{lemma}

\begin{proof}
By Lemmas~\ref{cb:lem:beqtok} and~\ref{cb:lem:beqlist}, $\mathsf{beqList}(\mathsf{beqTok},\cdot,\cdot)$ satisfies the
hypothesis of Lemma~\ref{cb:lem:beqlist} on the type of layers, and $\mathsf{beqTok}$ satisfies it on pairs of
naturals. Applying Lemma~\ref{cb:lem:beqlist} to the two conjuncts of $\mathsf{beqNode}(a,b)=\mathsf{true}$ gives
$a.\mathsf{layers}=b.\mathsf{layers}$ and $a.\mathsf{cube}=b.\mathsf{cube}$, hence $a=b$.
\end{proof}

\begin{lemma}[Converses; written step]\label{cb:lem:beq-refl}
$\mathsf{Nat.beq}(x,x)$, $\mathsf{beqTok}(t,t)$, $\mathsf{beqLayer}(L,L)$ and $\mathsf{beqNode}(n,n)$ are
$\mathsf{true}$ for every $x\in\N$, pair $t$, layer $L$ and node $n$. Consequently $\mathsf{beqTok}$,
$\mathsf{beqLayer}$, $\mathsf{beqNode}$ and $\mathsf{beqList}(\mathsf{Nat.beq},\cdot,\cdot)$ return $\mathsf{true}$
exactly on equal arguments.
\end{lemma}

\begin{proof}
If $eq(x,x)=\mathsf{true}$ for all $x$, then $\mathsf{beqList}(eq,l,l)=\mathsf{true}$ for every list $l$, by induction on
$l$. Apply this to $eq=\mathsf{Nat.beq}$ (reflexive), to $\mathsf{beqTok}$ (reflexive since $\mathsf{Nat.beq}$ is) and to
$\mathsf{beqList}(\mathsf{beqTok},\cdot,\cdot)$; for nodes both conjuncts are then $\mathsf{true}$. The other direction
is Lemmas~\ref{cb:lem:beqlist} (with $eq=\mathsf{Nat.beq}$ or $eq=\mathsf{beqTok}$), \ref{cb:lem:beqtok}
and~\ref{cb:lem:beqnode}.
\end{proof}

\subsection{Certificates and the checker}\label{cb:sec:checker}

\begin{definition}[Certificates; \leanref{C3Cover.Cert}{C3LEAN/C3Cover.lean:245}]\label{cb:def:cert}
A \emph{certificate} is a triple $c=\langle c.\mathsf{root},\ c.\mathsf{ds},\ c.\mathsf{ks}\rangle$ of natural numbers.
\end{definition}

\begin{definition}[Unpacking; \leanref{C3Cover.unpack}{C3LEAN/C3Cover.lean:251}]\label{cb:def:unpack}
For all $f,x\in\N$ we have $\mathsf{unpack}(0,x)=[]$; $\mathsf{unpack}(f+1,x)=[]$ if $x=0$; and
$\mathsf{unpack}(f+1,x)=(x\bmod2^{64})\mathbin{::}\mathsf{unpack}(f,\lfloor x/2^{64}\rfloor)$ if $x\ne0$ (in Lean
$2^{64}$ is written $18446744073709551616$).
\end{definition}

\begin{definition}[Keys of a certificate; \leanref{C3Cover.Cert.keys}{C3LEAN/C3Cover.lean:259}]\label{cb:def:keys}
For a certificate $c$, $\mathsf{keys}(c)=\mathsf{unpack}(100000,\ c.\mathsf{ks})$.
\end{definition}

\begin{definition}[The checker; \leanref{C3Cover.run}{C3LEAN/C3Cover.lean:262}]\label{cb:def:run}
For $f,ds\in\N$, a list $stk$ of nodes and a list $ks$ of naturals, the Boolean $\mathsf{run}(f,stk,ds,ks)$ is
defined by recursion on $f$ (the first applicable clause is used):
\begin{enumerate}[label=(R\arabic*)]
\item $\mathsf{run}(f,[],ds,ks)=(ds=0)\wedge\mathsf{isEmpty}(ks)$, for every $f$;
\item $\mathsf{run}(0,\ n\mathbin{::}rest,\ ds,ks)=\mathsf{false}$;
\item $\mathsf{run}(f+1,\ n\mathbin{::}rest,\ ds,ks)$ is computed from $d=ds\bmod16$ and $ds'=\lfloor ds/16\rfloor$:
  \begin{itemize}
  \item if $d=0$: it is $\mathsf{false}$ if $ks=[]$, and it is
    $\mathsf{beqNode}(n,\mathsf{decodeKey}(k))\wedge\mathsf{run}(f,rest,ds',ks')$ if $ks=k\mathbin{::}ks'$;
  \item if $d=1$: it is $\mathsf{decide}(2\le\card{n.\mathsf{layers}})\wedge
    \neg\mathsf{isEmpty}(\mathsf{layerChildren}(n))\wedge
    \mathsf{run}\big(f,\ \mathsf{layerChildren}(n)\mathbin{+\!\!+}rest,\ ds',\ ks\big)$;
  \item if $d\ge2$: it is $\mathsf{decide}(2\le\card{n.\mathsf{layers}})\wedge
    \mathsf{decide}(d\mathbin{\dot-}1\le\card{\mathsf{free}(n)})\wedge
    \mathsf{run}\big(f,\ \mathsf{sigmaChildren}(n,d\mathbin{\dot-}1)\mathbin{+\!\!+}rest,\ ds',\ ks\big)$.
  \end{itemize}
\end{enumerate}
\end{definition}

\begin{definition}[Root check; \leanref{C3Cover.checkRoot}{C3LEAN/C3Cover.lean:277}]\label{cb:def:checkroot}
For a node $\rho$ and a certificate $c$,
$\mathsf{checkRoot}(\rho,c)=\mathsf{run}(10^6,\ [\rho],\ c.\mathsf{ds},\ \mathsf{keys}(c))$ (in Lean $10^6$ is written
$1000000$).
\end{definition}

\begin{remark}[How a certificate is read]\label{cb:rem:cert}
The stack starts as $[\rho]$. At each step the first node $n$ is removed and the least significant base-$16$ digit $d$ of
$ds$ is consumed: $d=0$ declares $n$ a \emph{leaf}, which must be equal, as a node, to the decoded next key, and that
key is consumed; $d=1$ declares a layer split, and $d=1+K$ with $1\le K\le14$ a $\sigma$-split on the $K$ smallest
free residues of $n$ ($K$ $\sigma$ bits, $2^K$ children). For a
split the checker computes the children itself and puts them in front of the remaining stack. A certificate thus only
chooses the kind of split at each node and names the leaves, in depth-first order; it never supplies children. When
the stack is empty, the check succeeds iff no nonzero digit and no key is left. This makes the check strict (the
negative controls listed in \S\ref{cb:sec:coverage} exercise it), but it is not used in the soundness proof below.
Running out of fuel makes the check fail, so the bound $10^6$ is irrelevant for soundness; each step treats one node
of the cover tree, and the largest certificate has 849 of them (\dirn{C3LEAN/README.md}, \S1).
\end{remark}

\begin{lemma}[Checker soundness; \leanref{C3Cover.run_sound}{C3LEAN/C3Cover.lean:279}, \leanref{C3Cover.checkRoot_sound}{C3LEAN/C3Cover.lean:323}]\label{cb:lem:run}
The following hold.
\begin{enumerate}[label=(\alph*)]
\item (\lean{run_sound}) Let $S$ be a predicate on nodes and $f\in\N$. For every list $stk$ of nodes, every $ds\in\N$
  and every list $ks$ of naturals: if $\mathsf{run}(f,stk,ds,ks)=\mathsf{true}$ and $S(\mathsf{decodeKey}(k))$ holds
  for every $k\in ks$, then $\mathsf{Covered}(S)(n)$ holds for every $n\in stk$.
\item (\lean{checkRoot_sound}) Let $\rho$ be a node and $c$ a certificate with
  $\mathsf{checkRoot}(\rho,c)=\mathsf{true}$. Then $\mathsf{Covered}(S_c)(\rho)$, where $S_c(n)$ means: there is
  $k\in\mathsf{keys}(c)$ with $\mathsf{decodeKey}(k)=n$.
\end{enumerate}
\end{lemma}

\begin{proof}
We follow the Lean proof.

(a) Fix $S$. We prove by induction on $f$ the statement $\Phi(f)$: \emph{for all $stk$, $ds$, $ks$, if
$\mathsf{run}(f,stk,ds,ks)=\mathsf{true}$ and $S(\mathsf{decodeKey}(k))$ for every $k\in ks$, then
$\mathsf{Covered}(S)(n)$ for every $n\in stk$.} If $stk=[]$ there is nothing to show, for every $f$; so let
$stk=n_0\mathbin{::}rest$.

\emph{Case $f=0$.} By (R2), $\mathsf{run}(0,n_0\mathbin{::}rest,ds,ks)=\mathsf{false}$; the hypothesis fails, and
$\Phi(0)$ holds.

\emph{Step from $f$ to $f+1$.} Assume $\Phi(f)$, $\mathsf{run}(f+1,n_0\mathbin{::}rest,ds,ks)=\mathsf{true}$ and
$S(\mathsf{decodeKey}(k))$ for every $k\in ks$. Let $d=ds\bmod16$ and $ds'=\lfloor ds/16\rfloor$. We show that $n_0$
and every $n\in rest$ are covered, i.e.\ satisfy $\mathsf{Covered}(S)$.
\begin{itemize}
\item Case $d=0$. If $ks=[]$, the value would be $\mathsf{false}$; so $ks=k\mathbin{::}ks'$, with
  $\mathsf{beqNode}(n_0,\mathsf{decodeKey}(k))=\mathsf{true}$ and $\mathsf{run}(f,rest,ds',ks')=\mathsf{true}$. By
  Lemma~\ref{cb:lem:beqnode}, $n_0=\mathsf{decodeKey}(k)$; as $k\in ks$, $S(n_0)$ holds, and $n_0$ is covered by the
  rule (cert). Every $k'\in ks'$ lies in $ks$, so $S(\mathsf{decodeKey}(k'))$; hence $\Phi(f)$, applied to
  $(rest,ds',ks')$, shows that every $n\in rest$ is covered.
\item Case $d=1$. The three conjuncts are $\mathsf{true}$: $2\le\card{n_0.\mathsf{layers}}$;
  $\mathsf{isEmpty}(\mathsf{layerChildren}(n_0))=\mathsf{false}$, i.e.\ $\mathsf{layerChildren}(n_0)\ne[]$; and
  $\mathsf{run}(f,\mathsf{layerChildren}(n_0)\mathbin{+\!\!+}rest,ds',ks)=\mathsf{true}$. By $\Phi(f)$, applied to
  $(\mathsf{layerChildren}(n_0)\mathbin{+\!\!+}rest,ds',ks)$ with the same keys, every entry of
  $\mathsf{layerChildren}(n_0)$ and every entry of $rest$ is covered. The rule (layer) shows that $n_0$ is covered.
\item Case $d\ge2$. Put $K=d\mathbin{\dot-}1=d-1$, so $K\ge1>0$. The conjuncts give $2\le\card{n_0.\mathsf{layers}}$,
  $K\le\card{\mathsf{free}(n_0)}$ and $\mathsf{run}(f,\mathsf{sigmaChildren}(n_0,K)\mathbin{+\!\!+}rest,ds',ks)=
  \mathsf{true}$. By $\Phi(f)$ every entry of $\mathsf{sigmaChildren}(n_0,K)$ and of $rest$ is covered, and the rule
  (sigma) with this $K$ shows that $n_0$ is covered.
\end{itemize}
This proves $\Phi(f+1)$, and (a) follows. (Keys are consumed only at leaves; this is why the key hypothesis passes
unchanged to the recursive call in the two split cases.)

(b) Apply (a) with $S=S_c$, $f=10^6$, $stk=[\rho]$, $ds=c.\mathsf{ds}$ and $ks=\mathsf{keys}(c)$. The first hypothesis
of (a) is $\mathsf{checkRoot}(\rho,c)=\mathsf{true}$, and for every $k\in\mathsf{keys}(c)$ the statement
$S_c(\mathsf{decodeKey}(k))$ holds with witness $k$. Since $\rho\in[\rho]$, $\mathsf{Covered}(S_c)(\rho)$.
\end{proof}

\begin{lemma}[Concatenation; \leanref{C3Cover.all_append_of}{C3LEAN/C3Cover.lean:327}]\label{cb:lem:append}
Let $\alpha$ be a type, $f$ a Boolean-valued function on $\alpha$, and $l_1,l_2$ lists over $\alpha$. If
$\mathsf{all}(l_1,f)=\mathsf{true}$ and $\mathsf{all}(l_2,f)=\mathsf{true}$, then
$\mathsf{all}(l_1\mathbin{+\!\!+}l_2,f)=\mathsf{true}$.
\end{lemma}

\begin{proof}
By (L1), $\mathsf{all}(l_1\mathbin{+\!\!+}l_2,f)=\mathsf{all}(l_1,f)\wedge\mathsf{all}(l_2,f)=\mathsf{true}$. (Lean:
simplification with \lean{List.all_append}.)
\end{proof}

\subsection{Generated data and the computer-assisted facts}\label{cb:sec:data}

\begin{definition}[TSV root codes; \leanref{C3Cover.tsvRootCodes}{C3LEAN/C3Cover.lean:413}]\label{cb:def:tsv}
$\mathsf{tsvRootCodes}$ is an explicit list of natural numbers. According to its docstring and
\dirn{C3LEAN/README.md} (\S4) it is the root column of \dirn{C3FINAL/C3FINAL_mirrorh6_roots.tsv}, in file order,
written as root codes; this identification is checked by a Python script (\dirn{C3LEAN/tools/audit_keys.py}), not in
Lean.
\end{definition}

\begin{definition}[Certificate lists; \leanref{C3Cover.certs0}{C3LEAN/C3Cover.lean:423}]\label{cb:def:certsN}
$\mathsf{certs}_0,\dots,\mathsf{certs}_{126}$ (Lean \lean{certs0}, \dots, \lean{certs126}, at lines 423--1711; the
last, \lean{certs126}, at line 1668) are explicit lists of certificates, each certificate written as a triple of
numerals.
\end{definition}

\begin{definition}[All certificates; \leanref{C3Cover.certs}{C3LEAN/C3Cover.lean:1714}]\label{cb:def:certs}
The list of all certificates is, nested to the right,
\[
\mathsf{certs}=\mathsf{certs}_0\mathbin{+\!\!+}\big(\mathsf{certs}_1\mathbin{+\!\!+}\big(\cdots
\mathbin{+\!\!+}(\mathsf{certs}_{125}\mathbin{+\!\!+}\mathsf{certs}_{126})\cdots\big)\big).
\]
\end{definition}

\begin{definition}[Batch sizes; \leanref{C3Cover.bounds}{C3LEAN/C3Cover.lean:1717}]\label{cb:def:bounds}
$\mathsf{bounds}$ is an explicit list of 127 numbers, the sizes of the lists $\mathsf{certs}_N$ (between 1 and 85, with
sum 910). No theorem uses it.
\end{definition}

\begin{definition}[Key list; \leanref{C3Cover.certifiedKeys}{C3LEAN/C3Cover.lean:1737}]\label{cb:def:certifiedkeys}
$\mathsf{certifiedKeys}=\mathsf{flatMap}(\mathsf{keys},\mathsf{certs})$, the concatenation of the key lists of all
certificates. The \emph{leaf predicate} of the main theorem is
\[
S^\star(n)\quad:\Longleftrightarrow\quad\text{there is }k\in\mathsf{certifiedKeys}\text{ with }\mathsf{decodeKey}(k)=n.
\]
\end{definition}

The data were written by \dirn{C3LEAN/tools/gen_c3lean.py} from the run records of the project. Despite the name,
$\mathsf{certifiedKeys}$ is just the list of the keys of the certificates; nothing in Lean says which records exist.

\begin{definition}[Options; \leanref{C3Cover.isOption}{C3LEAN/C3Cover.lean:340}]\label{cb:def:isoption}
For a list $U$ of naturals and a layer $L$:
$\mathsf{isOption}(U,[(r,s)])=(r<11)\wedge\neg\mathsf{contains}(U,r)\wedge(s=2)$;
$\mathsf{isOption}(U,[(b,\sigma),(c,\tau)])=(b<c)\wedge(c<11)\wedge\neg\mathsf{contains}(U,b)\wedge
\neg\mathsf{contains}(U,c)\wedge(\sigma<2)\wedge(\tau<2)$; and $\mathsf{isOption}(U,L)=\mathsf{false}$ if $L$ has no
or more than two tokens.
\end{definition}

\begin{definition}[Distinct codes; \leanref{C3Cover.distinctCodes}{C3LEAN/C3Cover.lean:346}]\label{cb:def:distinctcodes}
For a list $l$ of naturals and $acc\in\N$: $\mathsf{distinctCodes}([],acc)=\mathsf{true}$ and
$\mathsf{distinctCodes}(x\mathbin{::}l,acc)=\neg\mathsf{testBit}(acc,x)\wedge\mathsf{distinctCodes}(l,\
acc\mathbin{\mathsf{or}}2^x)$.
\end{definition}

\begin{definition}[Code sets; \leanref{C3Cover.codeSet}{C3LEAN/C3Cover.lean:351}]\label{cb:def:codeset}
For a list $l$ of naturals, $\mathsf{codeSet}(l)=\mathsf{foldl}((acc,x)\mapsto acc\mathbin{\mathsf{or}}2^x,\ 0,\ l)$.
\end{definition}

\begin{definition}[Reflection facts; \leanref{C3Cover.reflectionFacts}{C3LEAN/C3Cover.lean:356}]\label{cb:def:reflectionfacts}
For a layer $t$, let $U=\mathsf{map}(u\mapsto u.1,t)$, $O=\mathsf{layerOptions}(U)$, $\iota(o)=\mathsf{refl}(t,o)$, and
let $\mathsf{kept}_t$ be as in Definition~\ref{cb:def:kept}. Then $\mathsf{reflectionFacts}(t)$ is the $6$-tuple of
the following values:
\begin{enumerate}[label=(F\arabic*)]
\item $\card{O}$;
\item $\mathsf{all}(O,\ \mathsf{isOption}(U,\cdot))\wedge\mathsf{distinctCodes}(\mathsf{map}(\mathsf{layerCode},O),\
  0)$;
\item $\mathsf{all}\big(O,\ o\mapsto\mathsf{isOption}(U,\iota(o))\wedge\mathsf{beqLayer}(\iota(\iota(o)),o)\big)$;
\item $\card{\mathsf{filter}(o\mapsto\mathsf{beqLayer}(\iota(o),o),\ O)}$;
\item $\card{\mathsf{keptOptions}(t)}$;
\item $\mathsf{all}\big(O,\ o\mapsto(\mathsf{kept}_t(o)\vee\mathsf{kept}_t(\iota(o)))\wedge
  (\mathsf{beqLayer}(\iota(o),o)\vee\neg(\mathsf{kept}_t(o)\wedge\mathsf{kept}_t(\iota(o))))\big)$.
\end{enumerate}
\end{definition}

\begin{lemma}[Distinctness and code sets; written step]\label{cb:lem:codes}
Let $l$ be a list of naturals and $acc\in\N$. Then $\mathsf{distinctCodes}(l,acc)=\mathsf{true}$ iff $l$ is
duplicate-free and no entry of $l$ lies in $\beta(acc)$. In particular $\mathsf{distinctCodes}(l,0)=\mathsf{true}$ iff $l$ is
duplicate-free. Moreover $\beta(\mathsf{codeSet}(l))=\mathrm{set}(l)$, so for lists $l,l'$ of naturals
$\mathsf{codeSet}(l)=\mathsf{codeSet}(l')$ iff $\mathrm{set}(l)=\mathrm{set}(l')$.
\end{lemma}

\begin{proof}
Induction on $l$, for all $acc$. For $l=[]$ both sides hold. For $l=x\mathbin{::}l'$ the value is $\mathsf{true}$ iff
$x\notin\beta(acc)$ and $\mathsf{distinctCodes}(l',acc\mathbin{\mathsf{or}}2^x)=\mathsf{true}$; by induction and
$\beta(acc\mathbin{\mathsf{or}}2^x)=\beta(acc)\cup\{x\}$, this holds iff $x\notin\beta(acc)$, $l'$ is duplicate-free,
and no entry of $l'$ lies in $\beta(acc)\cup\{x\}$; that is exactly the statement that $x\mathbin{::}l'$ is
duplicate-free with no entry in $\beta(acc)$. For $\beta(0)=\emptyset$ this gives the second sentence. The claim on
$\mathsf{codeSet}$ is Lemma~\ref{cb:lem:bits} with $z=0$ and $g$ the identity, together with the injectivity of
$\beta$.
\end{proof}

\begin{fact}[Batches; \leanref{C3Cover.batch0}{C3LEAN/C3Cover.lean:1740}]\label{cb:fact:batches}
(Computer-assisted.) For every $N\in\{0,\dots,126\}$,
\[
\mathsf{all}\big(\mathsf{certs}_N,\ c\mapsto\mathsf{checkRoot}(\mathsf{decodeRoot}(c.\mathsf{root}),c)\big)=
\mathsf{true}.
\]
These are the 127 theorems \lean{C3Cover.batch0}, \lean{C3Cover.batch1}, \dots,
\lean{C3Cover.batch126}; the one for $N$ is at line $1740+4N$ (the last, \lean{batch126}, at line 2244).
\end{fact}

Each of the 127 theorems is Lean-checked by kernel evaluation (\texttt{decide +kernel}); no hand
proof is given. What is evaluated: for each of the 910 certificates $c$ in the lists $\mathsf{certs}_N$, the node
$\mathsf{decodeRoot}(c.\mathsf{root})$ and the list $\mathsf{keys}(c)$ are computed, and $\mathsf{run}$ is evaluated
step by step as in Definition~\ref{cb:def:run}; at every split node the kernel computes $\mathsf{layerChildren}$, or
$\mathsf{free}$ and $\mathsf{sigmaChildren}$, from the definitions, and at every leaf it compares the node with the
decoded key by $\mathsf{beqNode}$. Over all 910 certificates this treats 76{,}256 nodes: 73{,}567 leaves, 1{,}747
$\sigma$-splits and 942 layer splits (counts recorded in \dirn{C3LEAN/c3cover_summary.json}; the review's independent
transcription of the checker into Python obtained the same counts, \dirn{C3LEAN/review/logs/indep_audit.json}). For
554 of the 910 roots the decision field $c.\mathsf{ds}$ of the certificate $c$ is the numeral $0$ and
$\mathsf{keys}(c)$ has one entry, so the certificate consists of the root itself as the only leaf. This number is our
own count from the certificate data of \dirn{C3LEAN/C3Cover.lean}; it is not a Lean statement, and no project record
states it.

\begin{lemma}[All checked; \leanref{C3Cover.all_checked}{C3LEAN/C3Cover.lean:2247}]\label{cb:lem:all-checked}
(Computer-assisted: the proof uses Fact~\ref{cb:fact:batches}.) Every certificate passes:
\[
\mathsf{all}\big(\mathsf{certs},\ c\mapsto\mathsf{checkRoot}(\mathsf{decodeRoot}(c.\mathsf{root}),c)\big)=\mathsf{true}.
\]
\end{lemma}

\begin{proof}
As in Lean. Let $p(c)=\mathsf{checkRoot}(\mathsf{decodeRoot}(c.\mathsf{root}),c)$, and let $T_{126}=\mathsf{certs}_{126}$
and $T_N=\mathsf{certs}_N\mathbin{+\!\!+}T_{N+1}$ for $N=125,\dots,0$, so that $T_0=\mathsf{certs}$ by
Definition~\ref{cb:def:certs}. By Fact~\ref{cb:fact:batches}, $\mathsf{all}(T_{126},p)=\mathsf{true}$; and if
$\mathsf{all}(T_{N+1},p)=\mathsf{true}$, then Lemma~\ref{cb:lem:append} together with Fact~\ref{cb:fact:batches} for
$\mathsf{certs}_N$ gives $\mathsf{all}(T_N,p)=\mathsf{true}$. Descending from $N=125$ to $N=0$ (126 applications of
\lean{all_append_of}, as in the Lean term) gives the claim.
\end{proof}

\begin{fact}[Reflection facts below H3; \leanref{C3Cover.reflection_facts_H3}{C3LEAN/C3Cover.lean:367}]\label{cb:fact:refl-h3}
(Computer-assisted in Lean, by kernel evaluation; written proof in Proposition~\ref{cb:prop:refl}.)
$\mathsf{reflectionFacts}(\mathsf{H3})=(190,\mathsf{true},\mathsf{true},10,100,\mathsf{true})$.
\end{fact}

\begin{fact}[Reflection facts below H6; \leanref{C3Cover.reflection_facts_H6_1_0}{C3LEAN/C3Cover.lean:370}]\label{cb:fact:refl-h6}
(Computer-assisted in Lean, by kernel evaluation; written proof in Proposition~\ref{cb:prop:refl}.)
For each $a\in\{1,\dots,5\}$ and $s\in\{0,1\}$:
$\mathsf{reflectionFacts}(\mathsf{H6}(a,s))=(153,\mathsf{true},\mathsf{true},9,81,\mathsf{true})$. These are the ten
theorems \lean{reflection_facts_H6_1_0}, \lean{reflection_facts_H6_1_1}, \lean{reflection_facts_H6_2_0}, \dots,
\lean{reflection_facts_H6_5_1}; the one for $(a,s)$ is at line $370+3(2a+s-2)$ (the last at line 397).
\end{fact}

\begin{fact}[Top cases; \leanref{C3Cover.topCases_eq}{C3LEAN/C3Cover.lean:401}]\label{cb:fact:topcases}
(Computer-assisted in Lean, by kernel evaluation; written proof in Proposition~\ref{cb:prop:refl}.)
The list of top cases is
\[
\mathsf{topCases}=[\mathsf{H3},\mathsf{H6}(1,0),\mathsf{H6}(1,1),\mathsf{H6}(2,0),\mathsf{H6}(2,1),\mathsf{H6}(3,0),
\mathsf{H6}(3,1),\mathsf{H6}(4,0),\mathsf{H6}(4,1),\mathsf{H6}(5,0),\mathsf{H6}(5,1)].
\]
\end{fact}

\begin{fact}[Strata; \leanref{C3Cover.strata}{C3LEAN/C3Cover.lean:406}]\label{cb:fact:strata}
(Computer-assisted in Lean, by kernel evaluation; written proof in Proposition~\ref{cb:prop:refl}.)
The numbers of kept options with one and with two tokens are
\begin{align*}
\mathsf{map}\big(t\mapsto\card{\mathsf{filter}(o\mapsto(\card o=1),\ \mathsf{keptOptions}(t))},\
\mathsf{topCases}\big)&=\mathsf{replicate}(11,5),\\
\mathsf{map}\big(t\mapsto\card{\mathsf{filter}(o\mapsto(\card o=2),\ \mathsf{keptOptions}(t))},\
\mathsf{topCases}\big)&=95\mathbin{::}\mathsf{replicate}(10,76).
\end{align*}
\end{fact}

Each of the 13 theorems of Facts~\ref{cb:fact:refl-h3}--\ref{cb:fact:strata} is Lean-checked by kernel evaluation of
the definitions; Proposition~\ref{cb:prop:refl} below gives a complete written proof of all of them. They are not used
in the proof of Theorem~\ref{cb:thm:c3}; they record that $\mathsf{roots}$ has the properties that the reductions of
the second layer need (Remark~\ref{cb:rem:use}).

\begin{fact}[Roots; \leanref{C3Cover.roots_match}{C3LEAN/C3Cover.lean:1720}]\label{cb:fact:roots-match}
(Computer-assisted in Lean, by kernel evaluation; written proof in Proposition~\ref{cb:prop:roots-match}.)
The list of roots has $910$ entries, and each is recovered from its code:
\[
\card{\mathsf{roots}}=910\quad\text{and}\quad
\mathsf{all}\big(\mathsf{roots},\ \rho\mapsto\mathsf{beqNode}(\mathsf{decodeRoot}(\mathsf{rootCode}(\rho)),\rho)\big)
=\mathsf{true}.
\]
\end{fact}

Lean-checked by kernel evaluation of the definitions (the kernel computes the list $\mathsf{roots}$);
Proposition~\ref{cb:prop:roots-match} below gives a complete written proof that does not use the computation.

\begin{fact}[Root codes; \leanref{C3Cover.roots_codes_match}{C3LEAN/C3Cover.lean:1725}]\label{cb:fact:roots-codes}
(Computer-assisted.) The codes of the roots match the roots of the certificates:
\[
\mathsf{beqList}\big(\mathsf{Nat.beq},\ \mathsf{map}(\mathsf{rootCode},\mathsf{roots}),\
\mathsf{map}(c\mapsto c.\mathsf{root},\mathsf{certs})\big)=\mathsf{true}.
\]
\end{fact}

\begin{fact}[TSV codes; \leanref{C3Cover.roots_eq_tsv}{C3LEAN/C3Cover.lean:1730}]\label{cb:fact:roots-tsv}
(Computer-assisted.) The following four statements hold (as one conjunction):
\begin{gather*}
\card{\mathsf{tsvRootCodes}}=910,\qquad \mathsf{distinctCodes}(\mathsf{tsvRootCodes},0)=\mathsf{true},\\
\mathsf{distinctCodes}(\mathsf{map}(c\mapsto c.\mathsf{root},\mathsf{certs}),0)=\mathsf{true},\\
\mathsf{codeSet}(\mathsf{map}(c\mapsto c.\mathsf{root},\mathsf{certs}))=\mathsf{codeSet}(\mathsf{tsvRootCodes}).
\end{gather*}
\end{fact}

Facts~\ref{cb:fact:roots-codes} and~\ref{cb:fact:roots-tsv} are computer-assisted and Lean-checked by kernel
evaluation: the kernel compares the 910 codes $\mathsf{rootCode}(\rho)$, computed from the definitions, entry by entry
and in order with the 910 generated numerals $c.\mathsf{root}$, and it tests, with the bit-set accumulator of
$\mathsf{distinctCodes}$ and the or-fold $\mathsf{codeSet}$, that the latter and the 910 numerals of
$\mathsf{tsvRootCodes}$ are duplicate-free lists with equal sets. No hand proof of them is given.

\subsection{Written proofs of the facts about the roots}\label{cb:sec:written}

Call a token $(r,s)$ \emph{valid} if $r\le10$ and $s\le2$, and a layer \emph{valid} if it consists of one or two valid
tokens.

\begin{lemma}[Sorting short lists; written step]\label{cb:lem:sort}
$\mathsf{sortTok}([t])=[t]$; $\mathsf{sortTok}([t,u])=[t,u]$ if $t.1\le u.1$, and $\mathsf{sortTok}([t,u])=[u,t]$ if
$t.1>u.1$. In particular a two-token list with distinct residues is put into increasing order of residue, and a layer
of one token, or of two tokens with increasing residues, is fixed by $\mathsf{sortTok}$.
\end{lemma}

\begin{proof}
$\mathsf{sortTok}([t])=\mathsf{insTok}(t,[])=[t]$, and $\mathsf{sortTok}([t,u])=\mathsf{insTok}(t,[u])$, which is
$[t,u]$ if $t.1\le u.1$ and $u\mathbin{::}\mathsf{insTok}(t,[])=[u,t]$ otherwise.
\end{proof}

\begin{lemma}[Layer options; written step]\label{cb:lem:options}
The following hold.
\begin{enumerate}[label=(\alph*)]
\item Let $l$ be a strictly increasing list of naturals. Then $\mathsf{pairOptions}(l)$ is duplicate-free, has
  $4\binom{\card l}{2}$ entries, and its entries are exactly the lists $[(b,\sigma),(c,\tau)]$ with $b,c\in l$, $b<c$
  and $\sigma,\tau\in\{0,1\}$.
\item Let $U$ be a list of naturals, $R=\{r\in\{0,\dots,10\}:r\notin U\}$ and $f=\card R$. Then
  $\mathsf{layerOptions}(U)$ is duplicate-free, has $f+4\binom f2$ entries, and its set of entries is
  \[
  \mathcal O(U)=\{[(r,2)]:r\in R\}\cup\{[(b,\sigma),(c,\tau)]:b,c\in R,\ b<c,\ \sigma,\tau\in\{0,1\}\}.
  \]
  Moreover, for every layer $L$: $\mathsf{isOption}(U,L)=\mathsf{true}$ iff $L\in\mathcal O(U)$.
\item Every element of $\mathcal O(U)$ is a valid layer and is fixed by $\mathsf{sortTok}$.
\end{enumerate}
\end{lemma}

\begin{proof}
(a) Induction on $l$. For $l=[]$ the list is empty and $\binom02=0$. Let $l=b\mathbin{::}l'$; then $l'$ is strictly
increasing and all its entries exceed $b$. The first part $\mathsf{flatMap}(c\mapsto Q_{b,c},l')$ consists of the
$4\card{l'}$ lists $[(b,\sigma),(c,\tau)]$ with $c\in l'$ and $\sigma,\tau\in\{0,1\}$; they are pairwise distinct
(distinct triples $(c,\sigma,\tau)$, as $l'$ is duplicate-free) and all have first residue $b$. By induction,
$\mathsf{pairOptions}(l')$ is duplicate-free with $4\binom{\card{l'}}2$ entries, all with first residue in $l'$, hence
different from $b$. So $\mathsf{pairOptions}(l)$ is duplicate-free with
$4\card{l'}+4\binom{\card{l'}}2=4\binom{\card{l'}+1}2$ entries, and each entry has the stated form ($b<c$ in the first
part since $c\in l'$; by induction in the second). Conversely, let $b',c'\in l$ with $b'<c'$ and
$\sigma,\tau\in\{0,1\}$. If $b'=b$, then $c'\ne b$, so $c'\in l'$, and $[(b,\sigma),(c',\tau)]$ is in the first part.
If $b'\ne b$, then $b'\in l'$, so $c'>b'>b$ and $c'\in l'$; by induction the list is in $\mathsf{pairOptions}(l')$.

(b) $\mathsf{rem}(U)$ is the list of the elements of $R$ in increasing order (filtering keeps the order of
$\mathsf{range}(11)$). The first part $\mathsf{map}(r\mapsto[(r,2)],\mathsf{rem}(U))$ is duplicate-free with $f$
one-token entries; the second part $\mathsf{pairOptions}(\mathsf{rem}(U))$ is described by (a), and its entries have
two tokens. So $\mathsf{layerOptions}(U)$ is duplicate-free with $f+4\binom f2$ entries, and its set of entries is
$\mathcal O(U)$. By Definition~\ref{cb:def:isoption}, $\mathsf{isOption}(U,[(r,s)])=\mathsf{true}$ iff $r\le10$,
$r\notin U$ and $s=2$; $\mathsf{isOption}(U,[(b,\sigma),(c,\tau)])=\mathsf{true}$ iff $b<c\le10$, $b,c\notin U$ and
$\sigma,\tau\in\{0,1\}$; and $\mathsf{isOption}(U,L)=\mathsf{false}$ for the other layers. This is membership in
$\mathcal O(U)$.

(c) The residues are at most $10$, the second components lie in $\{0,1,2\}$, and a two-token element has increasing
residues; apply Lemma~\ref{cb:lem:sort}.
\end{proof}

\begin{lemma}[Decoding; written step]\label{cb:lem:decode}
The following hold.
\begin{enumerate}[label=(\alph*)]
\item $\mathsf{decodeTok}(\mathsf{tokCode}(r,s))=(r,s)$ for all $r\in\N$ and $s\le2$.
\item For every valid layer $L$: $1\le\mathsf{layerCode}(L)\le1122$ and
  $\mathsf{decodeLayer}(\mathsf{layerCode}(L))=L$. In particular $\mathsf{layerCode}$ is injective on valid layers.
\item For valid layers $A,B$:
  $\mathsf{decodeRoot}\big(\mathsf{rootCode}(\langle[A,B],[]\rangle)\big)=\langle[A,B],[]\rangle$.
\end{enumerate}
\end{lemma}

\begin{proof}
(a) $\mathsf{tokCode}(r,s)=3r+s$ with $0\le s<3$, so $\lfloor(3r+s)/3\rfloor=r$ and $(3r+s)\bmod3=s$.

(b) A valid token $t=(r,s)$ has $\mathsf{tokCode}(t)=3r+s\le30+2=32$. If $L=[t]$, then
$c=\mathsf{layerCode}(L)=1+\mathsf{tokCode}(t)\in[1,33]$; as $c\le33$,
$\mathsf{decodeLayer}(c)=[\mathsf{decodeTok}(c\mathbin{\dot-}1)]=[\mathsf{decodeTok}(\mathsf{tokCode}(t))]=[t]$ by (a).
If $L=[t,u]$, let $x=\mathsf{tokCode}(t)$ and $y=\mathsf{tokCode}(u)$, both in $[0,32]$. Then
$c=34+33x+y\in[34,\ 34+33\cdot32+32]=[34,1122]$; as $c>33$, $\mathsf{decodeLayer}(c)$ is computed from
$c\mathbin{\dot-}34=33x+y$, and since $0\le y<33$, $\lfloor(33x+y)/33\rfloor=x$ and $(33x+y)\bmod33=y$; by (a),
$\mathsf{decodeLayer}(c)=[\mathsf{decodeTok}(x),\mathsf{decodeTok}(y)]=[t,u]$.

(c) $\mathsf{rootCode}(\langle[A,B],[]\rangle)=q\cdot1123+r$ with $q=\mathsf{layerCode}(A)$ and
$r=\mathsf{layerCode}(B)\in[1,1122]$ by (b). As $0\le r<1123$, $\lfloor(q\cdot1123+r)/1123\rfloor=q$ and
$(q\cdot1123+r)\bmod1123=r$. Hence $\mathsf{decodeRoot}$ returns
$\langle[\mathsf{decodeLayer}(q),\mathsf{decodeLayer}(r)],[]\rangle=\langle[A,B],[]\rangle$ by (b).
\end{proof}

\begin{lemma}[The order lexLe; written step]\label{cb:lem:lex}
For all lists $x,y$ of naturals: (a) $\mathsf{lexLe}(x,x)=\mathsf{true}$; (b) $\mathsf{lexLe}(x,y)=\mathsf{true}$ or
$\mathsf{lexLe}(y,x)=\mathsf{true}$; (c) if $\mathsf{lexLe}(x,y)=\mathsf{lexLe}(y,x)=\mathsf{true}$, then $x=y$.
\end{lemma}

\begin{proof}
Induction on $x$, for all $y$.
(a) $\mathsf{lexLe}([],[])=\mathsf{true}$, and
$\mathsf{lexLe}(a\mathbin{::}as,a\mathbin{::}as)=(a<a)\vee((a=a)\wedge\mathsf{lexLe}(as,as))=\mathsf{lexLe}(as,as)$,
which is $\mathsf{true}$ by induction.
(b) If $x=[]$ or $y=[]$, the empty list is $\mathsf{lexLe}$ the other one. Let $x=a\mathbin{::}as$ and
$y=b\mathbin{::}bs$. If $a<b$, then $\mathsf{lexLe}(x,y)=\mathsf{true}$; if $b<a$, then $\mathsf{lexLe}(y,x)=\mathsf{true}$;
if $a=b$, then $\mathsf{lexLe}(x,y)=\mathsf{lexLe}(as,bs)$ and $\mathsf{lexLe}(y,x)=\mathsf{lexLe}(bs,as)$, and one of
them is $\mathsf{true}$ by induction.
(c) If $x=[]$, then $\mathsf{lexLe}(y,[])=\mathsf{true}$ forces $y=[]$. If $x=a\mathbin{::}as$, then
$\mathsf{lexLe}(x,y)=\mathsf{true}$ forces $y=b\mathbin{::}bs$, and the hypotheses say: $a<b$ or ($a=b$ and
$\mathsf{lexLe}(as,bs)=\mathsf{true}$); and $b<a$ or ($b=a$ and $\mathsf{lexLe}(bs,as)=\mathsf{true}$). Since $a<b$
excludes both $b<a$ and $b=a$, and $b<a$ excludes $a=b$, we get $a=b$ and
$\mathsf{lexLe}(as,bs)=\mathsf{lexLe}(bs,as)=\mathsf{true}$; so $as=bs$ by induction, and $x=y$.
\end{proof}

\begin{lemma}[Counting kept options; written step]\label{cb:lem:orbits}
Let $X$ be a duplicate-free list of valid layers, each fixed by $\mathsf{sortTok}$, and let $\iota$ be a map from
layers to layers with $\iota(x)\in X$ and $\iota(\iota(x))=x$ for every $x\in X$. For $x\in X$ let
$k(x)=\mathsf{lexLe}(\mathsf{okey}(x),\mathsf{okey}(\iota(x)))$, and let $F$ be the number of $x\in X$ with
$\iota(x)=x$. Then:
\begin{enumerate}[label=(\alph*)]
\item if $\iota(x)=x$, then $k(x)=\mathsf{true}$;
\item if $\iota(x)\ne x$, then exactly one of $k(x)$, $k(\iota(x))$ is $\mathsf{true}$;
\item $\card{\mathsf{filter}(k,X)}=F+(\card X-F)/2$.
\end{enumerate}
\end{lemma}

\begin{proof}
(a) $k(x)=\mathsf{lexLe}(\mathsf{okey}(x),\mathsf{okey}(x))=\mathsf{true}$ by Lemma~\ref{cb:lem:lex}(a).

(b) For $y\in X$ we have $\mathsf{sortTok}(y)=y$, so $\mathsf{okey}(y)=\mathsf{map}(\kappa,y)$. On tokens with second
component $s\le2$, $\kappa(r,s)=3r+e(s)$ with $e(2)=0$, $e(0)=1$, $e(1)=2$, and $e$ is a bijection of $\{0,1,2\}$;
so $r=\lfloor\kappa(r,s)/3\rfloor$ and $e(s)=\kappa(r,s)\bmod3$ determine $(r,s)$, and $\kappa$ is injective on valid
tokens. Hence $\mathsf{map}(\kappa,\cdot)$ is injective on lists of valid tokens (equal images have equal lengths and
equal entries), and $\mathsf{okey}$ is injective on $\mathrm{set}(X)$. Let $\iota(x)\ne x$; both lie in $X$, so
$\mathsf{okey}(x)\ne\mathsf{okey}(\iota(x))$. Moreover
$k(\iota(x))=\mathsf{lexLe}(\mathsf{okey}(\iota(x)),\mathsf{okey}(\iota(\iota(x))))
=\mathsf{lexLe}(\mathsf{okey}(\iota(x)),\mathsf{okey}(x))$. By Lemma~\ref{cb:lem:lex}(b) at least one of $k(x)$,
$k(\iota(x))$ is $\mathsf{true}$, and by Lemma~\ref{cb:lem:lex}(c) not both.

(c) As $X$ is duplicate-free, $\card{\mathsf{filter}(k,X)}$ is the number of elements $x$ of $\mathrm{set}(X)$ with
$k(x)=\mathsf{true}$. The map $\iota$ restricted to $\mathrm{set}(X)$ is an involution of $\mathrm{set}(X)$, so
$\mathrm{set}(X)$ is the disjoint union of its $F$ fixed points and of $(\card X-F)/2$ two-element orbits
$\{x,\iota(x)\}$. By (a) every fixed point is counted, and by (b) exactly one element of each two-element orbit.
\end{proof}

\begin{proposition}[Proof of the reflection facts; written step]\label{cb:prop:refl}
Facts~\ref{cb:fact:refl-h3}, \ref{cb:fact:refl-h6}, \ref{cb:fact:topcases} and~\ref{cb:fact:strata} hold. More
precisely, let $t$ be $\mathsf{H3}$, with $R=\{1,\dots,10\}$, or $\mathsf{H6}(a,s)$ with $1\le a\le5$, $s\in\{0,1\}$,
with $R=\{1,\dots,10\}\setminus\{a\}$; let $O=\mathsf{layerOptions}(\mathsf{map}(u\mapsto u.1,t))$ and
$\iota=\mathsf{refl}(t,\cdot)$. Then $\iota$ maps $\mathrm{set}(O)$ into itself, is an involution of it, preserves the
number of tokens, and its fixed points are:
\begin{itemize}
\item for $\mathsf{H3}$: the ten layers $[(b,\sigma),(11-b,1-\sigma)]$ with $1\le b\le5$ and $\sigma\in\{0,1\}$;
\item for $\mathsf{H6}(a,s)$: the layer $[(r_0,2)]$, where $r_0=a/2$ if $a$ is even and $r_0=(a+11)/2$ if $a$ is odd
  (so $r_0=6,1,7,2,8$ for $a=1,\dots,5$), and the eight layers $[(b,\sigma),(b^*,\sigma\oplus f_b)]$ with $b\in R$,
  $b<b^*$ and $\sigma\in\{0,1\}$; here $r^*=((a+11)\mathbin{\dot-}r)\bmod11$ for $r\in R$ (that is, $r^*=a-r$ if $r<a$
  and $r^*=a+11-r$ if $r>a$), and $f_b=1$ if exactly one of $b<a$, $s=0$ holds and $f_b=0$ otherwise.
\end{itemize}
\end{proposition}

\begin{proof}
Lean checks these facts by evaluation; the following proof does not use the evaluation.

\emph{The used residues and the options.} For $t=\mathsf{H3}$, $\mathsf{map}(u\mapsto u.1,t)=[0]$; for
$t=\mathsf{H6}(a,s)$ it is $[0,a]$. So $R$ is the set of Lemma~\ref{cb:lem:options}(b), with $f=\card R=10$, resp.\
$9$. By that lemma, $O$ is duplicate-free with $\mathrm{set}(O)=\mathcal O$ (the set $\mathcal O(U)$ of the lemma),
and $\card O=10+4\binom{10}2=190$, resp.\ $9+4\binom92=153$. Let $X_1$ and $X_2$ be the sublists of one-token and of
two-token entries of $O$; they have $10$ and $180$, resp.\ $9$ and $144$, entries.

\emph{The map below $\mathsf{H3}$.} Here $\iota=\mathsf{mirror}$ (first case of Definition~\ref{cb:def:refl}). For
$r\in\{1,\dots,10\}$, $(11\mathbin{\dot-}r)\bmod11=11-r\in\{1,\dots,10\}$, and for $\sigma\in\{0,1\}$,
$1\mathbin{\dot-}\sigma=1-\sigma$. Hence, by Lemma~\ref{cb:lem:sort},
\[
\mathsf{mirror}([(r,2)])=[(11-r,2)],\qquad
\mathsf{mirror}([(b,\sigma),(c,\tau)])=[(11-c,1-\tau),(11-b,1-\sigma)]\quad(b<c),
\]
the second because $11-b>11-c$. Both results lie in $\mathcal O$, and applying the formulas twice returns the original
layer ($11-(11-r)=r$, $1-(1-\sigma)=\sigma$). So $\iota$ is an involution of $\mathcal O$ preserving the number of
tokens. A one-token option is never fixed ($11-r=r$ is impossible). A two-token option $[(b,\sigma),(c,\tau)]$, $b<c$,
is fixed iff $11-c=b$, $1-\tau=\sigma$, $11-b=c$ and $1-\sigma=\tau$, i.e.\ iff $c=11-b$ and $\tau=1-\sigma$; with
$b<c=11-b$ this means $1\le b\le5$. So $F=10$: none in $X_1$ and ten in $X_2$.

\emph{The map below $\mathsf{H6}(a,s)$.} Here $\iota=\mathsf{phi}(a,s,\cdot)$ (second case of
Definition~\ref{cb:def:refl}, with this $a$ and $s$). For $r\in R$ let $r^*=((a+11)\mathbin{\dot-}r)\bmod11$, as
in the statement. Since $r\le10<a+11$, $(a+11)\mathbin{\dot-}r=a+11-r$. If $r<a$, then $a+11-r\in[12,a+10]$, so $r^*=a-r\in[1,a-1]$; if $r>a$,
then $a+11-r\in[a+1,10]$, so $r^*=a+11-r\in[a+1,10]$. In both cases $r^*\in R$, $[r^*<a]=[r<a]$ and $(r^*)^*=r$ (indeed
$a-(a-r)=r$, resp.\ $a+11-(a+11-r)=r$). So $r\mapsto r^*$ is an involution of $R$. For $\sigma\in\{0,1\}$ the second
component of $\varphi_{a,s}(r,\sigma)$ is $\sigma\oplus f_r$, where $f_r=1$ if exactly one of $r<a$, $s=0$ holds and
$f_r=0$ otherwise (because $1\mathbin{\dot-}\sigma=\sigma\oplus1$); for $\sigma=2$ it is $2$. As $[r^*<a]=[r<a]$, we
have $f_{r^*}=f_r$. Hence, by Lemma~\ref{cb:lem:sort},
\[
\mathsf{phi}(a,s,[(r,2)])=[(r^*,2)],\qquad
\mathsf{phi}(a,s,[(b,\sigma),(c,\tau)])=\mathsf{sortTok}\big([(b^*,\sigma\oplus f_b),(c^*,\tau\oplus f_c)]\big),
\]
the latter being these two tokens in increasing order of residue ($b^*\ne c^*$ since $r\mapsto r^*$ is injective).
Both results lie in $\mathcal O$. The token map $h(r,\sigma)=(r^*,\sigma\oplus f_r)$, $h(r,2)=(r^*,2)$, satisfies
$h(h(x))=x$ on the tokens of options, since $(r^*)^*=r$ and $\sigma\oplus f_r\oplus f_{r^*}=\sigma$. For a two-token
option $o$, $\mathsf{sortTok}(\mathsf{map}(h,o))$ is $\mathsf{map}(h,o)$ or its reversal; applying $\mathsf{map}(h,\cdot)$
gives $o$ or its reversal, and $\mathsf{sortTok}$ of that is $o$, since $o$ has increasing residues. So $\iota$ is an
involution of $\mathcal O$ preserving the number of tokens.

Fixed points: $[(r,2)]$ is fixed iff $r^*=r$, i.e.\ iff $a=2r$ (case $r<a$) or $a+11=2r$ (case $r>a$). For even $a$
only the first equation has an integer solution, $r_0=a/2\in[1,a-1]$; for odd $a$ only the second,
$r_0=(a+11)/2\in[a+1,10]$ (as $a\le5$). So $r\mapsto r^*$ has exactly one fixed point $r_0$ in $R$, $X_1$ has exactly
one fixed entry, and $R\setminus\{r_0\}$ splits into four two-element orbits $\{b,b^*\}$. Let $o=[(b,\sigma),(c,\tau)]$
with $b<c$. If $\iota(o)=o$, the residues satisfy $\{b^*,c^*\}=\{b,c\}$; $b^*=b$ and $c^*=c$ is impossible (only one
fixed residue), so $c=b^*$ and $b=c^*$. In that case
$\iota(o)=\mathsf{sortTok}([(c,\sigma\oplus f_b),(b,\tau\oplus f_c)])=[(b,\tau\oplus f_c),(c,\sigma\oplus f_b)]$,
which equals $o$ iff $\tau\oplus f_c=\sigma$ and $\sigma\oplus f_b=\tau$; as $f_c=f_{b^*}=f_b$, both conditions say
$\tau=\sigma\oplus f_b$. So each of the four orbits $\{b<c\}$ carries exactly two fixed options (one for each
$\sigma$), and $F=9$: one in $X_1$ and eight in $X_2$.

\emph{Counting.} The lists $O$, $X_1$, $X_2$ satisfy the hypotheses of Lemma~\ref{cb:lem:orbits} with this $\iota$
(Lemma~\ref{cb:lem:options}(c); $\iota$ maps $\mathcal O$ to itself, preserves the number of tokens and is an
involution), and then $k=\mathsf{kept}_t$. Hence the numbers of kept entries are:
for $\mathsf{H3}$, $10+(190-10)/2=100$ in $O$, $0+10/2=5$ in $X_1$ and $10+(180-10)/2=95$ in $X_2$;
for $\mathsf{H6}(a,s)$, $9+(153-9)/2=81$ in $O$, $1+(9-1)/2=5$ in $X_1$ and $8+(144-8)/2=76$ in $X_2$.

\emph{The six components.} (F1) is $\card O=190$, resp.\ $153$. (F2): $\mathsf{all}(O,\mathsf{isOption}(U,\cdot))=
\mathsf{true}$ by Lemma~\ref{cb:lem:options}(b); the list $\mathsf{map}(\mathsf{layerCode},O)$ is duplicate-free
because $O$ is and $\mathsf{layerCode}$ is injective on valid layers (Lemma~\ref{cb:lem:decode}(b)), so
$\mathsf{distinctCodes}(\cdot,0)=\mathsf{true}$ by Lemma~\ref{cb:lem:codes}. (F3): for $o\in O$, $\iota(o)\in\mathcal O$,
so $\mathsf{isOption}(U,\iota(o))=\mathsf{true}$, and $\iota(\iota(o))=o$, so $\mathsf{beqLayer}(\iota(\iota(o)),o)=
\mathsf{true}$ by Lemma~\ref{cb:lem:beq-refl}. (F4): by Lemma~\ref{cb:lem:beq-refl},
$\mathsf{beqLayer}(\iota(o),o)=\mathsf{true}$ iff $\iota(o)=o$, so the value is $F=10$, resp.\ $9$. (F5):
$\mathsf{keptOptions}(t)=\mathsf{filter}(\mathsf{kept}_t,O)$ has $100$, resp.\ $81$, entries. (F6): for $o\in O$, by
Lemma~\ref{cb:lem:orbits}(a),(b) at least one of $\mathsf{kept}_t(o)$, $\mathsf{kept}_t(\iota(o))$ is $\mathsf{true}$;
and if $\mathsf{beqLayer}(\iota(o),o)=\mathsf{false}$, then $\iota(o)\ne o$ and not both are $\mathsf{true}$. This
proves Facts~\ref{cb:fact:refl-h3} and~\ref{cb:fact:refl-h6}.

\emph{Top cases and strata.} Unfolding Definition~\ref{cb:def:topcases} with $\mathsf{range}(5)=[0,1,2,3,4]$ gives the
list of Fact~\ref{cb:fact:topcases}. In $\mathsf{keptOptions}(t)$ the entries with one token are the kept entries of
$X_1$ and those with two tokens the kept entries of $X_2$. In the order of $\mathsf{topCases}$ ($\mathsf{H3}$ first,
then the ten $\mathsf{H6}(a,s)$), the counts are therefore $5$ eleven times, and $95$ followed by $76$ ten times. This
is Fact~\ref{cb:fact:strata}.
\end{proof}

\begin{proposition}[Proof of roots\_match; written step]\label{cb:prop:roots-match}
$\card{\mathsf{roots}}=910$, and for every $\rho\in\mathsf{roots}$,
$\mathsf{decodeRoot}(\mathsf{rootCode}(\rho))=\rho$ and
$\mathsf{beqNode}(\mathsf{decodeRoot}(\mathsf{rootCode}(\rho)),\rho)=\mathsf{true}$. Hence
Fact~\ref{cb:fact:roots-match} holds.
\end{proposition}

\begin{proof}
By (L3), $\card{\mathsf{roots}}=\sum_t\card{\mathsf{keptOptions}(t)}$ over the 11 entries of $\mathsf{topCases}$, which
is $100+10\cdot81=910$ by Proposition~\ref{cb:prop:refl}. By (L2) and (L4), every $\rho\in\mathsf{roots}$ is
$\langle[t,o],[]\rangle$ with $t\in\mathsf{topCases}$ and $o\in\mathsf{keptOptions}(t)$, so
$o\in\mathcal O(\mathsf{map}(u\mapsto u.1,t))$ (Lemma~\ref{cb:lem:options}(b)). The layer $t$ is valid
($\mathsf{H3}=[(0,0)]$, or $\mathsf{H6}(a,s)=[(0,0),(a,s)]$ with $a\le5$ and $s\le1$, by the unfolding of
$\mathsf{topCases}$ in Proposition~\ref{cb:prop:refl}), and $o$ is valid by Lemma~\ref{cb:lem:options}(c).
Lemma~\ref{cb:lem:decode}(c) gives $\mathsf{decodeRoot}(\mathsf{rootCode}(\rho))=\rho$, and
Lemma~\ref{cb:lem:beq-refl} gives $\mathsf{beqNode}(\rho,\rho)=\mathsf{true}$.
\end{proof}

\begin{corollary}[Roots and certificates; written step]\label{cb:cor:roots}
(Computer-assisted: the proof uses Facts~\ref{cb:fact:roots-codes} and~\ref{cb:fact:roots-tsv}.) The following hold.
\begin{enumerate}[label=(\alph*)]
\item $\mathsf{map}(\mathsf{rootCode},\mathsf{roots})=\mathsf{map}(c\mapsto c.\mathsf{root},\mathsf{certs})$ as lists,
  and $\card{\mathsf{certs}}=910$.
\item The 910 entries of $\mathsf{roots}$ are pairwise distinct nodes, and
  $\{\mathsf{rootCode}(\rho):\rho\in\mathsf{roots}\}=\mathrm{set}(\mathsf{tsvRootCodes})$.
\end{enumerate}
\end{corollary}

\begin{proof}
(a) Lemma~\ref{cb:lem:beqlist} with $eq=\mathsf{Nat.beq}$ (whose value $\mathsf{true}$ implies equality) turns
Fact~\ref{cb:fact:roots-codes} into the equality of lists; then
$\card{\mathsf{certs}}=\card{\mathsf{map}(\mathsf{rootCode},\mathsf{roots})}=910$ by (L3) and
Fact~\ref{cb:fact:roots-match}.
(b) By Fact~\ref{cb:fact:roots-tsv} and Lemma~\ref{cb:lem:codes}, $\mathsf{map}(c\mapsto c.\mathsf{root},\mathsf{certs})$
is duplicate-free, hence so is $\mathsf{map}(\mathsf{rootCode},\mathsf{roots})$ by (a); two positions of
$\mathsf{roots}$ holding the same node would give equal codes, so $\mathsf{roots}$ is duplicate-free. The equality of
the $\mathsf{codeSet}$ values and Lemma~\ref{cb:lem:codes} give
$\mathrm{set}(\mathsf{map}(c\mapsto c.\mathsf{root},\mathsf{certs}))=\mathrm{set}(\mathsf{tsvRootCodes})$, and (a)
identifies the left side with the set of root codes of the roots.
\end{proof}

The roots are distinct also by a hand argument (distinct top cases and distinct options give distinct nodes), but their
identification with $\mathsf{tsvRootCodes}$ rests on the computer-assisted Facts~\ref{cb:fact:roots-codes}
and~\ref{cb:fact:roots-tsv}.

\subsection{The threefold cover theorem}\label{cb:sec:c3}

\begin{theorem}[Roots covered; \leanref{C3Cover.all_roots_covered}{C3LEAN/C3Cover.lean:2251}]\label{cb:thm:c3}
(Computer-assisted: the proof uses Facts~\ref{cb:fact:batches} and~\ref{cb:fact:roots-codes}.) For every
$\rho\in\mathsf{roots}$, $\mathsf{Covered}(S^\star)(\rho)$, where $S^\star(n)$ means that there is
$k\in\mathsf{certifiedKeys}$ with $\mathsf{decodeKey}(k)=n$.
\end{theorem}

The theorem is Lean-checked (axioms \lean{propext}, \lean{Quot.sound}) and has no hypotheses. Its proof uses the
computer-assisted Facts~\ref{cb:fact:batches} (through Lemma~\ref{cb:lem:all-checked}) and~\ref{cb:fact:roots-codes},
and Fact~\ref{cb:fact:roots-match}, which also has the written proof of Proposition~\ref{cb:prop:roots-match}.

\begin{proof}
We follow the Lean proof. Let $\rho\in\mathsf{roots}$.
\begin{enumerate}[label=(\arabic*)]
\item By Fact~\ref{cb:fact:roots-match} (or Proposition~\ref{cb:prop:roots-match}), we have
  $\mathsf{beqNode}(\mathsf{decodeRoot}(\mathsf{rootCode}(\rho)),\rho)=\mathsf{true}$, so
  $\mathsf{decodeRoot}(\mathsf{rootCode}(\rho))=\rho$ by Lemma~\ref{cb:lem:beqnode}.
\item By Corollary~\ref{cb:cor:roots}(a), $\mathsf{map}(\mathsf{rootCode},\mathsf{roots})=\mathsf{map}(c\mapsto
  c.\mathsf{root},\mathsf{certs})$. As $\mathsf{rootCode}(\rho)$ is an entry of the left side (L4), there is
  $c\in\mathsf{certs}$ with $c.\mathsf{root}=\mathsf{rootCode}(\rho)$.
\item By Lemma~\ref{cb:lem:all-checked}, $\mathsf{checkRoot}(\mathsf{decodeRoot}(c.\mathsf{root}),c)=\mathsf{true}$.
  By (1) and (2), $\mathsf{decodeRoot}(c.\mathsf{root})=\mathsf{decodeRoot}(\mathsf{rootCode}(\rho))=\rho$, so
  $\mathsf{checkRoot}(\rho,c)=\mathsf{true}$.
\item By Lemma~\ref{cb:lem:run}(b), $\mathsf{Covered}(S_c)(\rho)$, with $S_c(n)$: there is $k\in\mathsf{keys}(c)$ with
  $\mathsf{decodeKey}(k)=n$.
\item If $S_c(n)$ holds with witness $k\in\mathsf{keys}(c)$, then $k\in\mathsf{flatMap}(\mathsf{keys},\mathsf{certs})=
  \mathsf{certifiedKeys}$ by (L2), since $c\in\mathsf{certs}$; so $S^\star(n)$. Lemma~\ref{cb:lem:mono} turns
  $\mathsf{Covered}(S_c)(\rho)$ into $\mathsf{Covered}(S^\star)(\rho)$.\qedhere
\end{enumerate}
\end{proof}

\begin{remark}[What the theorem says, and how the project uses it]\label{cb:rem:use}
\emph{The leaf set.} Theorem~\ref{cb:thm:c3} is only as informative as the predicate $S^\star$, which is read off the
certificates themselves. A certificate that names its root as the only leaf satisfies the checker; this is the case for
554 of the 910 roots (our count from the certificate data; see the paragraph after Fact~\ref{cb:fact:batches}), and
for those roots the statement is immediate. The substantial
claim of the project, that every node satisfying $S^\star$ names a run record in which both checkers are recorded as
accepting the proof (the criterion of \dirn{C3U2/tools/compose.py}; \dirn{C3LEAN/README.md}, \S4) and whose CNF hash
matches the base formula plus the units of the node, is checked outside Lean by
the Python scripts \dirn{C3LEAN/tools/audit_keys.py} and \dirn{C3LEAN/review/indep_audit.py} (73{,}567 leaves,
73{,}707 records, no mismatch; \dirn{C3LEAN/review/logs/indep_audit.json}). The review demonstrates the gap with a
certificate that collapses an inner $\sigma$-split of root~0 into one leaf naming a node without such a record and
still passes (\lean{C3Cover.rev_uncertified_leaf_passes}, \dirn{C3LEAN/review/RevCtrl_pass.lean:307}). The verdicts
themselves are recorded fields, and the proofs were not kept (\dirn{C3LEAN/README.md}, \S4). The project later recorded that
acceptance by \texttt{lrat-check} alone is not a certificate and that \texttt{lrat-check} verdicts are supporting only
(lines 111--114 of \dirn{WRITEUP/note/symmetric_extremal.tex}; \dirn{LRATCHECK20261105}); of the two recorded
verdicts, only the \texttt{cake\_lpr} verdict counts as a certificate (\dirn{WRITEUP/SYMMETRIC.md}, \S3).

\emph{Intended use (not established here).} By the induction principle of Definition~\ref{cb:def:covered}, the theorem
gives $P(\rho)$ for every root $\rho$ and every $S^\star$-closed predicate $P$. In the project, $P(n)$ is ``no set of
the class (threefold symmetric, 33 points, general position, no convex $7$-gon, with its canonical labelling) is
contained in the node $n$''. That $P$ is $S^\star$-closed amounts to three claims: (i) $P(n)$ for every leaf, which is
the record-based fact above together with the soundness of the CNF encoding; (ii) the rule (layer), which is Lemma~4
of \dirn{C3-S/soundness.md} and, for nodes with a nonempty cube, rule~3 of \dirn{C3U2/sigma_split.md} (Lean's rule
(layer) applies to such nodes, whereas the project's composer \dirn{C3U2/tools/compose.py} also requires a layer-split
record there, \dirn{C3LEAN/README.md}, \S5, finding~3); (iii) the rule (sigma), which is rule~2 of
\dirn{C3U2/sigma_split.md}. The passage from ``$P$ holds at every root'' to ``the class is empty'' is the cover
theorem with the reduced cut in \dirn{C3H6M/soundness_h6m.md}, which uses Lemma~5 of \dirn{C3-S/soundness.md} and
Lemma~H6M; the $90$ options below $\mathsf{H3}$ and the $720$ options below the ten $\mathsf{H6}(a,s)$ that these
reductions drop ($190-100$ and $10\cdot(153-81)$) are not roots of the file, so Theorem~\ref{cb:thm:c3} says nothing
about them. Facts~\ref{cb:fact:refl-h3}
and~\ref{cb:fact:refl-h6} record that the computed second layers have the properties those reductions need:
$\mathsf{refl}$ is an involution of the options, and the kept options contain exactly one member of each pair and every
fixed option. None of (i)--(iii), and none of the written arguments cited here, is Lean-checked or proved in this
document. The project's own list of what the Lean check leaves open, with the status it records for each part, is
\S6 of \dirn{C3LEAN/README.md}; its trust base for the threefold theorem (Theorem~1 of
\dirn{WRITEUP/SYMMETRIC.md}) is \S3 of that file.
\end{remark}

\subsection{The central 34-point split}\label{cb:sec:c2}

Every Lean reference in a statement head of this subsection is to \dirn{C2FINAL/Central34Orbits.lean}, namespace
\lean{Central34}; the statements tagged ``written step'' have no Lean counterpart.

\begin{definition}[Inner classes; \leanref{Central34.inner}{C2FINAL/Central34Orbits.lean:21}]\label{cb:def:inner}
For $a\in\N$,
\[
\mathsf{inner}(a)=\mathsf{filter}\big(i\mapsto(i\ne0)\wedge(i\ne a),\ \mathsf{range}(17)\big).
\]
\end{definition}

\begin{definition}[Pairs; \leanref{Central34.choose2}{C2FINAL/Central34Orbits.lean:24}]\label{cb:def:choose2}
For lists of naturals,
\[
\mathsf{choose2}([])=[],\qquad
\mathsf{choose2}(x\mathbin{::}xs)=\mathsf{map}\big(y\mapsto[x,y],\ xs\big)\mathbin{+\!\!+}\mathsf{choose2}(xs).
\]
\end{definition}

\begin{definition}[Triples; \leanref{Central34.choose3}{C2FINAL/Central34Orbits.lean:29}]\label{cb:def:choose3}
For lists of naturals,
\[
\mathsf{choose3}([])=[],\qquad
\mathsf{choose3}(x\mathbin{::}xs)=\mathsf{map}\big(W\mapsto x\mathbin{::}W,\ \mathsf{choose2}(xs)\big)
\mathbin{+\!\!+}\mathsf{choose3}(xs).
\]
\end{definition}

\begin{definition}[Subcases; \leanref{Central34.subcases}{C2FINAL/Central34Orbits.lean:34}]\label{cb:def:subcases}
For $a\in\N$, $\mathsf{subcases}(a)=\mathsf{choose2}(\mathsf{inner}(a))\mathbin{+\!\!+}\mathsf{choose3}(\mathsf{inner}(a))$.
\end{definition}

\begin{definition}[Masks; \leanref{Central34.mask}{C2FINAL/Central34Orbits.lean:37}]\label{cb:def:mask}
For a list $V$ of naturals,
\[
\mathsf{mask}(V)=\mathsf{foldl}\big((acc,i)\mapsto acc\mathbin{\mathsf{or}}2^i,\ 0,\ V\big).
\]
\end{definition}

\begin{definition}[Central reflection; \leanref{Central34.refl}{C2FINAL/Central34Orbits.lean:40}]\label{cb:def:c2refl}
For $a\in\N$ and a list $V$ of naturals,
\[
\mathsf{refl}(a,V)=\mathsf{map}\big(i\mapsto((a+17)\mathbin{\dot-}i)\bmod17,\ V\big).
\]
(This is \lean{Central34.refl}, not to be confused with \lean{C3Cover.refl}.)
\end{definition}

\begin{definition}[Bit sets; \leanref{Central34.bitset}{C2FINAL/Central34Orbits.lean:43}]\label{cb:def:bitset}
For a list $L$ of lists of naturals,
\[
\mathsf{bitset}(L)=\mathsf{foldl}\big((acc,W)\mapsto acc\mathbin{\mathsf{or}}2^{\mathsf{mask}(W)},\ 0,\ L\big).
\]
\end{definition}

\begin{definition}[Uncovered subcases; \leanref{Central34.uncovered}{C2FINAL/Central34Orbits.lean:46}]\label{cb:def:uncovered}
For $a\in\N$ and a list $C$ of lists of naturals, with the local abbreviation $B=\mathsf{bitset}(C)$,
\[
\mathsf{uncovered}(a,C)=\mathsf{filter}\Big(V\mapsto\neg\big(\mathsf{testBit}(B,\mathsf{mask}(V))\vee
\mathsf{testBit}(B,\mathsf{mask}(\mathsf{refl}(a,V)))\big),\ \mathsf{subcases}(a)\Big).
\]
\end{definition}

\begin{definition}[Closed lists; \leanref{Central34.closedP1}{C2FINAL/Central34Orbits.lean:51}]\label{cb:def:closedP}
The lists $\mathsf{closedP}_1,\dots,\mathsf{closedP}_8$ (Lean \lean{closedP1}, \dots, \lean{closedP8}, at lines 51,
78, 105, 132, 159, 186, 213, 240; the data end at line 264) are explicit lists of lists of naturals with $287$ entries
each, except $288$ for $\mathsf{closedP}_6$.
\end{definition}

\begin{definition}[Closed subcases; \leanref{Central34.closed}{C2FINAL/Central34Orbits.lean:267}]\label{cb:def:closed}
For $a\in\N$: $\mathsf{closed}(a)=\mathsf{closedP}_a$ if $1\le a\le8$, and $\mathsf{closed}(a)=[]$ otherwise.
\end{definition}

By Lemma~\ref{cb:lem:bits}, $\beta(\mathsf{mask}(V))=\mathrm{set}(V)$, so $\mathsf{mask}(V)=\mathsf{mask}(W)$ iff
$\mathrm{set}(V)=\mathrm{set}(W)$; and $\beta(\mathsf{bitset}(L))=\{\mathsf{mask}(W):W\in L\}$, so
$\mathsf{testBit}(\mathsf{bitset}(L),\mathsf{mask}(V))=\mathsf{true}$ iff $\mathrm{set}(V)=\mathrm{set}(W)$ for some
$W\in L$. For $1\le a\le16$ put $I_a=\{1,\dots,16\}\setminus\{a\}$.

\begin{remark}[Intended reading]\label{cb:rem:c2-reading}
In the project (header of the file and \dirn{C2FINAL/README.md}, \S1), a centrally symmetric $34$-point set consists of
$17$ antipodal classes $0,\dots,16$; in the hull case $P_a$, $1\le a\le8$, the hull is the parallelogram on the classes
$0$ and $a$, and the second convex layer is spanned by $2$ or $3$ of the other $15$ classes, which gives the subcases;
the reflection $i\mapsto(a-i)\bmod17$ is reduction~B. The lists $\mathsf{closedP}_a$ were generated by
\dirn{C2FINAL/gen_lean.py} as the subcases with a run record in which both checkers, \texttt{cake\_lpr} and
\texttt{lrat-check}, are recorded as accepting the proof and whose CNF hash was
rebuilt and matched by \dirn{C2FINAL/hashmatch.py} (or a re-solve of \dirn{C2FINAL/README.md}, \S4). This reading is not
used in any proof.
\end{remark}

\begin{lemma}[The subcases; written step]\label{cb:lem:c2-subsets}
Let $l$ be a strictly increasing list of $k$ naturals. Then $\mathsf{choose2}(l)$ is duplicate-free with $\binom k2$
entries, which are exactly the lists $[x,y]$ with $x,y\in l$, $x<y$; and $\mathsf{choose3}(l)$ is duplicate-free with
$\binom k3$ entries, which are exactly the lists $[x,y,z]$ with $x,y,z\in l$, $x<y<z$. Consequently, for
$1\le a\le16$, $\mathsf{subcases}(a)$ is a duplicate-free list of $\binom{15}2+\binom{15}3=105+455=560$ strictly
increasing lists, and $V\mapsto\mathrm{set}(V)$ is a bijection from its entries onto the set of subsets of $I_a$ with
two or three elements.
\end{lemma}

\begin{proof}
Induction on $l$; for $l=[]$ both lists are empty. Let $l=x\mathbin{::}xs$, so $xs$ is strictly increasing with $k-1$
entries, all larger than $x$. The first part of $\mathsf{choose2}(l)$ consists of the $k-1$ distinct lists $[x,y]$,
$y\in xs$, all increasing and with head $x$; by induction $\mathsf{choose2}(xs)$ is duplicate-free with
$\binom{k-1}2$ increasing entries, with heads in $xs$, hence different from $x$. So $\mathsf{choose2}(l)$ is
duplicate-free with $k-1+\binom{k-1}2=\binom k2$ entries. Every increasing pair $[u,v]$ of elements of $l$ occurs: if
$u=x$, then $v\in xs$ and $[x,v]$ is in the first part; otherwise $u,v\in xs$ (as $v>u>x$) and the pair is in
$\mathsf{choose2}(xs)$ by induction. Likewise the first part of $\mathsf{choose3}(l)$ consists of the
$\binom{k-1}2$ distinct lists $x\mathbin{::}W$ with $W\in\mathsf{choose2}(xs)$, increasing with head $x$, and
$\mathsf{choose3}(xs)$ contributes $\binom{k-1}3$ increasing triples with heads in $xs$; the total is
$\binom{k-1}2+\binom{k-1}3=\binom k3$, without repetition, and every increasing triple of elements of $l$ occurs (the
same case distinction on its first element). For $1\le a\le16$, $\mathsf{inner}(a)$ is the increasing list of the
$15$ elements of $I_a$. The pairs and the triples are disjoint by length, and a strictly increasing list is determined
by its set of entries; this gives the bijection.
\end{proof}

\begin{theorem}[Orbits; \leanref{Central34.subcases_reflection}{C2FINAL/Central34Orbits.lean:273}]\label{cb:thm:c2-refl}
(Computer-assisted in Lean, by kernel evaluation; written proof below.) The statement is
\[
\mathsf{all}\big(\mathsf{range}(8),\ k\mapsto Q(k+1)\big)=\mathsf{true},
\]
where for $a\in\N$, with $S_a=\mathsf{bitset}(\mathsf{subcases}(a))$, $Q(a)$ is the Boolean conjunction of
\begin{enumerate}[label=(\roman*)]
\item $\card{\mathsf{subcases}(a)}=560$;
\item $\mathsf{all}\big(\mathsf{subcases}(a),\ V\mapsto\mathrm{Img}_a(V)\wedge\mathrm{Inv}_a(V)\big)$, where
  $\mathrm{Img}_a(V)=\mathsf{testBit}(S_a,\mathsf{mask}(\mathsf{refl}(a,V)))$ and
  $\mathrm{Inv}_a(V)=\big(\mathsf{mask}(\mathsf{refl}(a,\mathsf{refl}(a,V)))=\mathsf{mask}(V)\big)$;
\item $\card{\mathsf{filter}\big(V\mapsto\mathsf{mask}(V)\le\mathsf{mask}(\mathsf{refl}(a,V)),\
  \mathsf{subcases}(a)\big)}=287$;
\item $\card{\mathsf{filter}\big(V\mapsto\mathsf{mask}(V)=\mathsf{mask}(\mathsf{refl}(a,V)),\
  \mathsf{subcases}(a)\big)}=14$.
\end{enumerate}
In words: for every $a\in\{1,\dots,8\}$, the statements (i)--(iv) hold.
\end{theorem}

The theorem is Lean-checked according to the record (paragraph ``Axioms'' in \S\ref{cb:sec:overview}), by kernel
evaluation of the definitions. The following written proof does not use the
evaluation.

\begin{proof}
Since $\mathsf{range}(8)=[0,\dots,7]$, the statement says that $Q(a)=\mathsf{true}$ for $a=1,\dots,8$. Fix such an
$a$ and let $\rho_a(i)=((a+17)\mathbin{\dot-}i)\bmod17$ for $i\in\{0,\dots,16\}$. As $i<a+17$,
$\rho_a(i)=(a+17-i)\bmod17$, so $\rho_a(i)\equiv a-i\pmod{17}$ and $\rho_a(i)\in\{0,\dots,16\}$; hence
$\rho_a(\rho_a(i))\equiv i$, and $\rho_a(\rho_a(i))=i$. Also $\rho_a(0)=a$ and $\rho_a(a)=0$. So $\rho_a$ is an
involution of $\{0,\dots,16\}$ that maps $I_a$ onto itself. For $V\in\mathsf{subcases}(a)$ the entries of $V$ lie in
$I_a$, $\mathsf{refl}(a,V)=\mathsf{map}(\rho_a,V)$ and $\mathrm{set}(\mathsf{refl}(a,V))=\rho_a(\mathrm{set}(V))$.

By Lemma~\ref{cb:lem:c2-subsets} there is, for each $V\in\mathsf{subcases}(a)$, a unique entry $V'$ of
$\mathsf{subcases}(a)$ with $\mathrm{set}(V')=\rho_a(\mathrm{set}(V))$, since $\rho_a(\mathrm{set}(V))\subseteq I_a$
has $\card{\mathrm{set}(V)}\in\{2,3\}$ elements; and $(V')'=V$ because $\rho_a\circ\rho_a$ is the identity. Moreover
$\mathsf{mask}(\mathsf{refl}(a,V))=\mathsf{mask}(V')$.

(i) is Lemma~\ref{cb:lem:c2-subsets}. (ii): $\mathrm{Img}_a(V)=\mathsf{true}$, because the number
$\mathsf{mask}(\mathsf{refl}(a,V))=\mathsf{mask}(V')$ lies in $\beta(S_a)$, as $V'\in\mathsf{subcases}(a)$; and
$\mathrm{Inv}_a(V)=\mathsf{true}$, because
$\mathrm{set}(\mathsf{refl}(a,\mathsf{refl}(a,V)))=\rho_a(\rho_a(\mathrm{set}(V)))=\mathrm{set}(V)$, so the two masks
agree.

(iv): $\mathsf{mask}(V)=\mathsf{mask}(\mathsf{refl}(a,V))$ iff $\mathrm{set}(V)=\rho_a(\mathrm{set}(V))$, i.e.\ iff
$\mathrm{set}(V)$ is a union of orbits of $\rho_a$ on $I_a$. A point $i\in I_a$ is fixed iff $2i\equiv a\pmod{17}$; as
$2\cdot9\equiv1\pmod{17}$, this means $i\equiv9a\pmod{17}$, so there is exactly one fixed point $x_0=9a\bmod17$ in
$\{0,\dots,16\}$, and $x_0\in I_a$ because $9a\not\equiv0$ and $9a\not\equiv a\pmod{17}$ for $1\le a\le8$ ($x_0=9,1,10,
2,11,3,12,4$ for $a=1,\dots,8$). The other $14$ elements of $I_a$ form $7$ two-element orbits. A union of orbits with
two elements is one of these $7$ orbits; one with three elements is $\{x_0\}$ together with one of them. So exactly
$7+7=14$ entries $V$ satisfy the condition.

(iii): if $\mathsf{mask}(V)=\mathsf{mask}(\mathsf{refl}(a,V))$, the condition $\le$ holds; this accounts for $14$
entries. The other $546$ entries are paired by $V\leftrightarrow V'$ with $V'\ne V$, and $\mathsf{mask}(V)$,
$\mathsf{mask}(V')$ are distinct naturals; as $\mathsf{mask}(\mathsf{refl}(a,V))=\mathsf{mask}(V')$ and
$\mathsf{mask}(\mathsf{refl}(a,V'))=\mathsf{mask}(V)$, exactly one member of each pair satisfies the condition. So the
count is $14+546/2=287$.
\end{proof}

\begin{fact}[Lists; \leanref{Central34.closed_are_subcases}{C2FINAL/Central34Orbits.lean:284}]\label{cb:fact:c2-sub}
(Computer-assisted.) The statement is
\[
\mathsf{all}\Big(\mathsf{range}(8),\ k\mapsto\mathsf{all}\big(\mathsf{closed}(k+1),\ V\mapsto
\mathsf{testBit}(\mathsf{bitset}(\mathsf{subcases}(k+1)),\mathsf{mask}(V))\big)\Big)=\mathsf{true}.
\]
Equivalently (by Lemma~\ref{cb:lem:c2-subsets} and the remarks after Definition~\ref{cb:def:closed}): for every
$a\in\{1,\dots,8\}$ and every $V\in\mathsf{closedP}_a$, $\mathrm{set}(V)$ is a subset of $I_a$ with two or three
elements.
\end{fact}

Lean-checked according to the record (paragraph ``Axioms'' in \S\ref{cb:sec:overview}), by kernel evaluation over
the explicit lists: for each $a$, the kernel computes the
masks of the $560$ subcases and their or-fold, and tests the mask of each of the $287$ (for $a=6$: $288$) entries of
$\mathsf{closedP}_a$ against it. No hand proof is given.

\begin{theorem}[Closure; \leanref{Central34.all_cases_closed}{C2FINAL/Central34Orbits.lean:291}]\label{cb:thm:c2}
(Computer-assisted.) The statement is
\[
\mathsf{all}\big(\mathsf{range}(8),\ k\mapsto\mathsf{isEmpty}(\mathsf{uncovered}(k+1,\mathsf{closed}(k+1)))\big)=
\mathsf{true};
\]
that is, $\mathsf{uncovered}(a,\mathsf{closedP}_a)=[]$ for every $a\in\{1,\dots,8\}$.
\end{theorem}

The theorem is Lean-checked according to the record (paragraph ``Axioms'' in \S\ref{cb:sec:overview}), by kernel
evaluation; no hand proof is given, and none is possible
without the explicit lists. What is evaluated: for each $a\in\{1,\dots,8\}$ the kernel computes
$B_a=\mathsf{bitset}(\mathsf{closedP}_a)$, the or-fold of $2^{\mathsf{mask}(W)}$ over the $287$ (for $a=6$: $288$)
entries $W$, and then, for each of the $560$ entries $V$ of $\mathsf{subcases}(a)$, the disjunction of the bits of
$B_a$ at $\mathsf{mask}(V)$ and at $\mathsf{mask}(\mathsf{refl}(a,V))$ (the second bit is needed only when the first
is $0$); the disjunction is $\mathsf{true}$ for every $V$, so the filtered list is empty. The two formulations in the statement agree because $\mathsf{range}(8)=[0,\dots,7]$,
$\mathsf{closed}(a)=\mathsf{closedP}_a$ for $1\le a\le8$, and $\mathsf{isEmpty}(l)=\mathsf{true}$ iff $l=[]$. The
statement is not vacuous: $\mathsf{uncovered}(a,[])=\mathsf{subcases}(a)$ has $560$ entries.

\begin{corollary}[Set form; written step]\label{cb:cor:c2}
(Computer-assisted: a consequence of Theorem~\ref{cb:thm:c2}.) Let $a\in\{1,\dots,8\}$, and for $i\in\{0,\dots,16\}$
let $\rho_a(i)=((a+17)\mathbin{\dot-}i)\bmod17$, the residue of $a-i$ modulo $17$. For every subset $T\subseteq I_a$
with two or three elements there is $W\in\mathsf{closedP}_a$ with $\mathrm{set}(W)=T$ or
$\mathrm{set}(W)=\rho_a(T)$.
\end{corollary}

\begin{proof}
This is a consequence of Theorem~\ref{cb:thm:c2}. By Lemma~\ref{cb:lem:c2-subsets}, $T=\mathrm{set}(V)$ for some
$V\in\mathsf{subcases}(a)$. Since $V$ is not an entry of $\mathsf{uncovered}(a,\mathsf{closedP}_a)=[]$,
$\mathsf{testBit}(B_a,\mathsf{mask}(V))$ or $\mathsf{testBit}(B_a,\mathsf{mask}(\mathsf{refl}(a,V)))$ is
$\mathsf{true}$, where $B_a=\mathsf{bitset}(\mathsf{closedP}_a)$. By the remarks after Definition~\ref{cb:def:closed},
there is $W\in\mathsf{closedP}_a$ with $\mathrm{set}(W)=\mathrm{set}(V)=T$ or
$\mathrm{set}(W)=\mathrm{set}(\mathsf{refl}(a,V))$; and $\mathrm{set}(\mathsf{refl}(a,V))=\rho_a(\mathrm{set}(V))=
\rho_a(T)$, because $\mathsf{refl}(a,V)=\mathsf{map}(\rho_a,V)$ by Definition~\ref{cb:def:c2refl}, the entries of $V$
lying in $I_a\subseteq\{0,\dots,16\}$.
\end{proof}

\begin{remark}[Scope of the central theorem]\label{cb:rem:c2-scope}
Theorem~\ref{cb:thm:c2} and Corollary~\ref{cb:cor:c2} say that the explicit lists meet every orbit $\{T,\rho_a(T)\}$ of
second-layer subcases, for each of the eight hull cases. That the listed subcases are refuted is not part of the
statement: the lists were selected outside Lean from run records (Remark~\ref{cb:rem:c2-reading}), and according to
\dirn{C2FINAL/README.md} (Status and \S9) most parallelogram pieces rest on recorded checker verdicts whose proofs
were deleted after checking. As in Remark~\ref{cb:rem:use}, of the two recorded checker verdicts only the
\texttt{cake\_lpr} verdict counts as a certificate; the project's trust base for its central theorem (Theorem~4 of
\dirn{WRITEUP/SYMMETRIC.md}) is \S3 of that file. The second-layer covers of the cases $P_a$, the $120$ hexagonal hull cases, the hull
cover, the reductions A ($P_a\sim P_{17-a}$) and B (the reflection), and the soundness of the encoder (\dirn{C2ENC})
are written arguments or records outside Lean and outside this document. By inspection of the data, the list
$\mathsf{closedP}_6$ contains both members of exactly one orbit, $[1,7]$ and $[5,16]$ (as sets,
$\{5,16\}=\rho_6(\{1,7\})$); this accounts for its $288$ entries. The header comment
of the file mentions a list \lean{pilotUnhashed} that the file does not define (stale text of the generator); it has
no effect on the theorems.
\end{remark}

\subsection{Coverage}\label{cb:sec:coverage}

\paragraph{Covered.} From \dirn{C3LEAN/C3Cover.lean}, every definition of the scope is covered by a definition above
(the generated family \lean{certs0}, \dots, \lean{certs126} by a single one):
\lean{m}, \lean{Tok}, \lean{Layer}, \lean{Node}, \lean{Node.used}, \lean{pairOptions}, \lean{layerOptions},
\lean{insTok}, \lean{sortTok}, \lean{topCases}, \lean{okey}, \lean{lexLe}, \lean{mirror}, \lean{phi}, \lean{refl},
\lean{keptOptions}, \lean{roots} (Definitions~\ref{cb:def:m}--\ref{cb:def:roots}); \lean{Node.free}, \lean{bitsOf},
\lean{addBits}, \lean{Node.sigmaChildren}, \lean{conflicts}, \lean{Node.layerChildren}
(Definitions~\ref{cb:def:free}--\ref{cb:def:layerchildren}); \lean{Covered} (Definition~\ref{cb:def:covered});
\lean{decodeTok}, \lean{decodeLayer}, \lean{decodeLayers}, \lean{decodeCube}, \lean{decodeKey}, \lean{tokCode},
\lean{layerCode}, \lean{rootCode}, \lean{decodeRoot}, \lean{beqList}, \lean{beqTok}, \lean{beqNode}, \lean{beqLayer}
(Definitions~\ref{cb:def:decodetok}--\ref{cb:def:beqlayer}); \lean{Cert}, \lean{unpack}, \lean{Cert.keys}, \lean{run},
\lean{checkRoot} (Definitions~\ref{cb:def:cert}--\ref{cb:def:checkroot}); \lean{tsvRootCodes},
\lean{certs0}--\lean{certs126}, \lean{certs}, \lean{bounds}, \lean{certifiedKeys}
(Definitions~\ref{cb:def:tsv}--\ref{cb:def:certifiedkeys}; the data are described, not reproduced); \lean{isOption},
\lean{distinctCodes}, \lean{codeSet}, \lean{reflectionFacts} (Definitions~\ref{cb:def:isoption}--%
\ref{cb:def:reflectionfacts}). The theorems: \lean{Covered.mono} (Lemma~\ref{cb:lem:mono}); \lean{beqList_eq},
\lean{beqTok_eq}, \lean{beqNode_eq} (Lemmas~\ref{cb:lem:beqlist}--\ref{cb:lem:beqnode}); \lean{run_sound} and
\lean{checkRoot_sound} (Lemma~\ref{cb:lem:run}); \lean{all_append_of} (Lemma~\ref{cb:lem:append});
\lean{batch0}--\lean{batch126} (Fact~\ref{cb:fact:batches}); \lean{all_checked} (Lemma~\ref{cb:lem:all-checked});
\lean{reflection_facts_H3} (Fact~\ref{cb:fact:refl-h3}); the ten \lean{reflection_facts_H6_1_0}, \dots,
\lean{reflection_facts_H6_5_1} (Fact~\ref{cb:fact:refl-h6}); \lean{topCases_eq} (Fact~\ref{cb:fact:topcases});
\lean{strata} (Fact~\ref{cb:fact:strata}); these four with the written proof of Proposition~\ref{cb:prop:refl};
\lean{roots_match} (Fact~\ref{cb:fact:roots-match}, with Proposition~\ref{cb:prop:roots-match});
\lean{roots_codes_match} and \lean{roots_eq_tsv} (Facts~\ref{cb:fact:roots-codes} and~\ref{cb:fact:roots-tsv}, with
Corollary~\ref{cb:cor:roots}); \lean{all_roots_covered} (Theorem~\ref{cb:thm:c3}). From
\dirn{C2FINAL/Central34Orbits.lean}: \lean{inner}, \lean{choose2}, \lean{choose3}, \lean{subcases}, \lean{mask},
\lean{refl}, \lean{bitset}, \lean{uncovered}, \lean{closedP1}--\lean{closedP8} (by a single definition), \lean{closed}
(Definitions~\ref{cb:def:inner}--\ref{cb:def:closed}); \lean{subcases_reflection} (Theorem~\ref{cb:thm:c2-refl}),
\lean{closed_are_subcases} (Fact~\ref{cb:fact:c2-sub}) and \lean{all_cases_closed} (Theorem~\ref{cb:thm:c2}, with
Corollary~\ref{cb:cor:c2}).

\paragraph{Not covered, deliberately.}
\begin{itemize}
\item The seven negative controls at the end of \dirn{C3LEAN/C3Cover.lean}, which are tests, not results:
  \lean{control_missing_record} (line 2266; root~0 with its fourth key removed is rejected),
  \lean{control_missing_sigma_child} (2270; root~0 without the subtree of the last child of its first $\sigma$-split),
  \lean{control_swapped_keys} (2274; two adjacent keys swapped), \lean{control_missing_layer_child} (2278; root~100,
  code $41560$, without the key of its last layer child), \lean{control_wrong_root} (2282; the certificate of root~0
  does not pass for root~1), \lean{control_strict_reduction} (2286; a strict comparison in the reduction would leave
  $810$ roots; this agrees with Proposition~\ref{cb:prop:refl}, as $810=(100-10)+10\cdot(81-9)$) and
  \lean{control_phi_mutant} (2292; $\mathsf{phi}$ without its $[s=0]$ term would keep options below $\mathsf{H6}(1,0)$
  whose root codes are not all in $\mathsf{tsvRootCodes}$). The \texttt{\#print axioms} commands (lines 2301--2327 of
  \dirn{C3LEAN/C3Cover.lean}, lines 296--298 of \dirn{C2FINAL/Central34Orbits.lean}) are not statements.
\item \dirn{C3LEAN/tools/prelude.lean} (SHA-256 \texttt{cdcf96de7739}\dots) is the verbatim source of lines 29--329 of
  \dirn{C3LEAN/C3Cover.lean}; its declarations are covered through that file.
\item \dirn{C3LEAN/Control_record.lean} (SHA-256 \texttt{5f755e897c30}\dots) and \dirn{C3LEAN/Control_sibling.lean}
  (\texttt{fb68b330e6e3}\dots): must-fail files, each a verbatim copy of the prelude plus one theorem (line 305)
  asserting that a mutated certificate of root~0 passes; they fail to compile, as intended
  (\dirn{C3LEAN/logs/control_record.log}, \dirn{C3LEAN/logs/control_sibling.log}).
\item \dirn{C3LEAN/review/RevCtrl_mut.lean} (\texttt{49bf2cce1b73}\dots): a copy of the prelude plus nine reviewer
  controls stating that mutated certificates are rejected (\lean{rev_extra_decision} has the same certificate as
  \lean{rev_extra_key}). \dirn{C3LEAN/review/RevCtrl_pass.lean} (\texttt{a135e2b92e6a}\dots): a copy of the prelude plus
  two controls stating that two non-genuine certificates pass (\lean{rev_uncertified_leaf_passes}, used in
  Remark~\ref{cb:rem:use}; \lean{rev_noncanonical_key_passes}) and the evaluation \lean{rev_parent_decodes} (used in
  Remark~\ref{cb:rem:keys}). \dirn{C3LEAN/review/RevDiff.lean} (\texttt{5aa3f6172026}\dots): a copy of the prelude plus
  an \texttt{\#eval} harness, without theorems of its own. These review files are tests and audits, not results;
  their prelude copies duplicate covered declarations.
\item The Lean package \dirn{C3L}: \dirn{C3L/C3Encoding.lean} (only imports) with \dirn{C3L/C3Encoding/Gen.lean},
  \dirn{C3L/C3Encoding/Case.lean} and \dirn{C3L/C3Encoding/Emit.lean}. It is a Lean program that writes, for
  parameters $m$ and $k$, the base CNF of the project's threefold encoding and the node CNFs (the base CNF followed by
  the unit clauses of a node) as DIMACS files, and it states no theorem (\dirn{C3L/README.md}). Its $41$ declarations
  are $39$ definitions, a structure and an abbreviation. In \dirn{C3L/C3Encoding/Gen.lean}, namespace \lean{C3}: the
  enumerations \lean{pairs}, \lean{triples}, \lean{quads}, \lean{perms2}, \lean{others}; the helpers \lean{key3},
  \lean{paritySort}, \lean{rot}, \lean{eps}; the variables and literal tables \lean{orbits}, the structure
  \lean{Tables}, \lean{mkTables}, \lean{Tables.s}, \lean{Tables.sig}, \lean{Tables.Cl}, \lean{Tables.Tl},
  \lean{Tables.invv}, \lean{Tables.Bvv}, \lean{Tables.B}, \lean{Tables.ccw}, \lean{chainVars}, \lean{numVars}; the
  clauses \lean{addCl}, \lean{push}, \lean{xor0}, \lean{secSig}, \lean{secSb}, \lean{secPrefix}, \lean{secForced},
  \lean{secXorB}, \lean{secXorT}, \lean{secChain} and the base formula \lean{baseClauses}. In
  \dirn{C3L/C3Encoding/Case.lean}, namespace \lean{C3}: the abbreviation \lean{Layer}, \lean{parseLayers},
  \lean{layerOrder} and \lean{layerUnits} (the unit clauses of a node). In \dirn{C3L/C3Encoding/Emit.lean}, root
  namespace: \lean{clauseLine}, \lean{body}, \lean{writeCnf} and \lean{main}, the output routines of the executable
  \texttt{c3\_dimacs}. That the output agrees with the CNF files of the project's Python generator is recorded in
  \dirn{C3L/logs} (comparisons outside Lean); the soundness theorem planned in \dirn{C3L/PLAN.md} (named
  \texttt{c3\_base\_sound} there) is not in the repository. CNF encodings are outside this document (item~(ii) of the
  paragraph ``What is not established here'' in \S\ref{cb:sec:overview}).
\end{itemize}

\paragraph{Treated in another section.} The Lean-checked files \dirn{AUDIT3/A1_reduction/SizeArith.lean} and
\dirn{AUDIT3/A1_reduction/verify_A1-F1/AssembledSizes.lean} (the size arithmetic of the reduction of the symmetric
$k=7$ statement to its case classes) and \dirn{ES7SPLIT/review/equivalence/BaekSeven.lean} (an instance of Baek's
theorem for abstract configurations) are outside the scope of this section; they are treated in the last section,
on size arithmetic.

\paragraph{Audit scripts of results outside Stage~1.} The five Lean files of \dirn{AUDIT2}, an adversarial review of
the project's mirror-class result, state no theorem of their own. \dirn{AUDIT2/A1/A1Dump.lean} defines \lean{cube}
(line~5) and \lean{main} (line~6), an input--output program that writes the data \lean{V2L.caseTable} (with the
leaves of the trees \lean{V2L.treeOf}), \lean{V2L.upTable} and \lean{V1L.hullLeaves} to a file, for comparison with a
Python decoding of their source files. \dirn{AUDIT2/A1/A1Audit.lean}, \dirn{AUDIT2/A2/Ax.lean},
\dirn{AUDIT2/A2/Ax2.lean} and \dirn{AUDIT2/A6/scripts/A6Check.lean} contain, apart from \texttt{import}, \texttt{open}
and \texttt{set\_option} lines and comments, only \texttt{\#check}, \texttt{\#print}, \texttt{\#print axioms} and
\texttt{\#eval} commands. The theorems they audit, such as \lean{V2L.mirrorExclusionLine_final},
\lean{V1L.hullVertexLine_of_leaves} and \lean{Mir.mirrorExclusionLine_of_unsat}, belong to the Lean developments in
\dirn{H3}, \dirn{V1-L} and \dirn{V2-L}, which lie outside Stage~1; this document does not restate them.
